% संपूर्ण मराठी ग्रंथातील प्रकरण 13: संपूर्णता प्रमेय
% हा संपादनयोग्य स्रोतखंड आहे. संकलनासाठी संपूर्ण स्रोतसंग्रहातील
% अक्षररूपे, पूर्वभाग, चिन्हव्याख्या आणि आकृती-स्रोत आवश्यक आहेत.
% मूळ: Open Logic Project; मजकूर आणि रूपांतर CC BY 4.0.
% यंत्रानुवाद, दुरुस्ती आणि तपासणी: OpenAI Codex —
% GPT-5.6 Sol आणि GPT-6 Sol, दोन्ही Ultra effort.
% दुरुस्ती-पुनरावलोकन: GPT-6.1 Sol, Ultra effort.
\chapter{संपूर्णता प्रमेय}\label{fol:com::chap}









\section{परिचय}\label{fol:com:int:sec}

संपूर्णता प्रमेय हा तर्कशास्त्राविषयीच्या सर्वांत मूलभूत निष्कर्षांपैकी एक
आहे. त्याच्या दोन मांडण्या आहेत; त्यांची सममूल्यता आपण सिद्ध करू. त्याची
पहिली मांडणी चिन्हार्थी परिणाम आणि आपल्या निष्पत्ती-पद्धतीतील संबंधाबद्दल
एक मूलभूत विधान करते: जर वाक्य~$\varphi$ हे काही वाक्ये
$\Gamma$ यांच्यापासून निष्पन्न होत असेल, तर $\Gamma \Proves \varphi$ हे
प्रस्थापित करणारी निष्पत्ती सुद्धा असते. त्यामुळे प्रत्यक्षात निष्पन्न
न होणाऱ्या गोष्टी सिद्ध न करता निष्पत्ती-पद्धत जितकी सामर्थ्यवान असू
शकते तितकी ती असते.

दुसऱ्या मांडणीत ते प्रतिमानाच्या अस्तित्वाविषयीचा निष्कर्ष म्हणून सांगता येते:
वाक्यांचा प्रत्येक सुसंगत संच पूर्ततायोग्य असतो. सुसंगतता ही
सिद्धता-उपपत्तीय संकल्पना आहे: आपली निष्पत्ती-पद्धत विशिष्ट
निष्पत्त्या निर्माण करू शकत नाही, असे ती सांगते. पण~$\Gamma$ पासून
विशिष्ट प्रकारच्या निष्पत्त्यांनाहीत एवढ्यावरूनच
संरचना~$\Struct{M}$
अस्तित्वात असून $\Sat{M}{\Gamma}$ आहे, याची हमी
आहे, असे कोण म्हणेल? संपूर्णता प्रमेय प्रथम सिद्ध होण्यापूर्वी---
खरे तर आज आपल्याकडे असलेल्या निष्पत्ती-पद्धती निर्माण होण्यापूर्वी---
महान जर्मन गणितज्ञ डाव्हीट हिल्बर्ट यांचे मत होते की गणिती सिद्धांतांची
सुसंगतता त्यांच्या विषयातील वस्तूंच्या अस्तित्वाची हमी देते. गोटलोप फ्रेग
यांना लिहिलेल्या पत्रात त्यांनी ते पुढीलप्रमाणे मांडले:
\begin{quote}
  स्वेच्छपणे दिलेली स्वयंसिद्धके आणि त्यांचे सर्व परिणाम परस्परविरोधी नसतील,
  तर ती सत्य असतात आणि स्वयंसिद्धकांनी परिभाषित केलेल्या वस्तू अस्तित्वात
  असतात. माझ्यासाठी सत्य आणि अस्तित्वाचा निकष हाच आहे.
\end{quote}
फ्रेग यांनी याला तीव्र विरोध केला. संपूर्णता प्रमेयाची दुसरी मांडणी दाखवते
की निदान पुढील अर्थाने हिल्बर्ट बरोबर होते: स्वयंसिद्धके सुसंगत असतील, तर
त्या सर्वांना सत्य करणारी \emph{एखादी}
संरचना अस्तित्वात असते.

संपूर्णता प्रमेय---किंवा अधिक नेमके सांगायचे तर त्याची सिद्धता---महत्त्वाची
असण्याची हीच काही एकमेव कारणे नाहीत. त्याचे अनेक महत्त्वाचे परिणाम आहेत;
त्यांपैकी काहींची स्वतंत्र चर्चा करू. उदाहरणार्थ, $\Gamma \Proves \varphi$ हे
दाखवणारी कोणतीही निष्पत्ती सांत असल्यामुळे ती~$\Gamma$ मधील सांतच
वाक्ये वापरू शकते. म्हणून संपूर्णता प्रमेयावरून असे निष्पन्न होते की
$\varphi$ हा~$\Gamma$ चा परिणाम असेल, तर तो~$\Gamma$ च्या एखाद्या सांत उपसंचाचाही
परिणाम असतो. याला \emph{संहतता} म्हणतात. सममूल्य रीतीने, $\Gamma$ चा प्रत्येक
सांत उपसंच सुसंगत असेल, तर $\Gamma$ स्वतः सुसंगत असला पाहिजे.

संहतता प्रमेय हे संपूर्णता प्रमेयापासून निष्पत्त्या मधून जाणाऱ्या वळणाने
मिळत असले, तरी संपूर्णता प्रमेयाची \emph{सिद्धता} वापरून ते थेट प्रस्थापित
करणेही शक्य आहे. कारण त्या सिद्धतेत विशिष्ट गुणधर्म---सुसंगतता---असलेला
वाक्यांचा संच घेऊन त्यातून विशिष्ट गुणधर्म असलेली संरचना
रचली जाते (या प्रकरणात ती त्या संचाची पूर्ति करते). सुसंगत संचांऐवजी
``सांततः पूर्ततायोग्य'' वाक्यांच्या संचांपासून सुरुवात केल्यास जवळजवळ
हीच रचना वापरून संहतता थेट प्रस्थापित करता येते.
या रचनेतून इतर परिणामही मिळतात; उदाहरणार्थ, वाक्यांच्या
  कोणत्याही पूर्ततायोग्य संचाला सांत किंवा गणनीय अनंत असे प्रतिमान असते.
  (या निष्कर्षाला ल्योव्हेनहाइम--स्कोलेम प्रमेय म्हणतात.) सर्वसाधारणपणे,
  वाक्यांच्या संचांपासून संरचना रचण्याची पद्धत तर्कशास्त्रात
  वारंवार आणि कधी कधी तत्त्वज्ञानातही वापरली जाते.












\section{सिद्धतेचा आराखडा}\label{fol:com:out:sec}

संपूर्णता प्रमेयाची सिद्धता काहीशी गुंतागुंतीची आहे आणि ती प्रथम वाचताना
वाट चुकणे सोपे आहे. म्हणून सिद्धतेचा आराखडा पाहू. पहिली पायरी म्हणजे
दृष्टिकोनातील बदल; त्यामुळे सिद्धतेकडे जाणारा मार्ग दिसू लागतो. संपूर्णतेचा
अर्थ “जेव्हा जेव्हा $\Gamma \Entails \varphi$, तेव्हा $\Gamma \Proves \varphi$”
असा घेतल्यास कल्पनाही सुचणे कठीण जाऊ शकते: कारण $\Gamma \Proves \varphi$ हे
दाखवण्यासाठी आपल्याला निष्पत्ती शोधावी लागते, आणि $\Gamma \Entails
\varphi$ हे गृहीतक त्यासाठी कोणत्याही प्रकारे मदत करते असे दिसत नाही. काही
सिद्धता-पद्धतींसाठी निष्पत्ती थेट रचता येते; पण आपण थोडा वेगळा मार्ग
घेऊ. आवश्यक दृष्टिकोन-बदल असा आहे: संपूर्णता पुढीलप्रमाणेही मांडता येते:
“$\Gamma$ सुसंगत असेल, तर तो पूर्ततायोग्य असतो.” कदाचित~$\Gamma$ मधील
माहिती आणि तो सुसंगत असल्याचे गृहीतक वापरून आपण~$\Gamma$ मधील प्रत्येक
वाक्य पूर्ण करणारी
संरचना रचू शकू. शेवटी, आपल्याला
कोणत्या प्रकारची संरचना हवी आहे हे
माहीत आहे:~$\Gamma$ जिचे वर्णन करतो अशीच!{}

$\Gamma$ मध्ये फक्त आण्विक
  वाक्ये असतील, तर त्याच्यासाठी
प्रतिमान रचणे सोपे आहे. सर्व आण्विक वाक्ये ही
  $\Atom{P}{a_1,\dots,a_n}$ या रूपाची आहेत असे समजा; येथे $a_i$ ही
  स्थिरांक आहेत. आपल्याला फक्त

  प्रांत~$\Domain{M}$ आणि~$P$ साठी असे निर्देशन ठरवायचे आहे की
  $\Sat{M}{\Atom{P}{a_1,\ldots,a_n}}$. पण ते फार कठीण नाही:
  $\Domain{M} = \Nat$, $\Assign{\Obj c_i}{M} = i$ असे ठेवा आणि प्रत्येक
  $\Atom{P}{a_1,\ldots,a_n} \in \Gamma$ साठी $\tuple{k_1,
    \dots, k_n}$ हे क्रमित बहुपद $\Assign{P}{M}$ मध्ये घाला; येथे $k_i$
  हा स्थिर चिन्ह~$a_i$ चा निर्देशांक आहे (म्हणजे $a_i \ident \Obj
  c_{k_i}$).

आता $\Gamma$ मध्ये~$\lnot \psi$ हे एखादे सूत्र असून $\psi$ आण्विक आहे
असे समजा. $\Struct{M}$ ची रचना $\lnot \psi$ सत्य
करण्याच्या शक्यतेत अडथळा आणेल, अशी काळजी वाटू शकते. पण येथेच~$\Gamma$ ची
सुसंगतता उपयोगी पडते: $\lnot \psi \in \Gamma$ असेल, तर $\psi \notin \Gamma$;
अन्यथा $\Gamma$ विसंगत असेल. आणि $\psi \notin \Gamma$ असेल, तर आपल्या
$\Struct{M}$ च्या रचनेनुसार
$\Sat/{M}{\psi}$; म्हणून
$\Sat{M}{\lnot \psi}$. आतापर्यंत सर्व ठीक आहे.

$\Gamma$ मध्ये गुंतागुंतीची, अनाण्विक सूत्रे असतील तर काय? उदाहरणार्थ,
त्यात $\varphi \land \psi$ आहे असे समजा. ते सत्य करण्यासाठी $\varphi$ आणि $\psi$ ही
दोन्ही~$\Gamma$ मध्ये असल्याप्रमाणे आपण पुढे गेले पाहिजे. आणि $\varphi \lor
\psi \in \Gamma$ असेल, तर त्यांपैकी किमान एक सत्य करावे लागेल; म्हणजे
त्यांपैकी एक~$\Gamma$ मध्ये असल्याप्रमाणे पुढे जावे लागेल.

यावरून पुढील कल्पना सुचते: आपण~$\Gamma$ मध्ये अतिरिक्त सूत्रे
अशा रीतीने भर घालतो की (a)~मिळणारा संच सुसंगत राहील, (b)~प्रत्येक
संभाव्य आण्विक वाक्य~$\varphi$ साठी मिळणाऱ्या संचात $\varphi$ किंवा
$\lnot \varphi$ यांपैकी एक असेल, आणि (c)~संचात $\varphi \land \psi$ असेल तेव्हा
$\varphi$ व $\psi$ ही दोन्ही असतील, तर $\varphi \lor \psi$ असेल तेव्हा $\varphi$ किंवा
$\psi$ यांपैकी किमान एकही असेल, इत्यादी. हे आपण करत राहतो (कदाचित
अनंतकाळपर्यंत). अशा रीतीने भर घातलेल्या सर्व सूत्रांच्या संचाला
$\Gamma^*$ म्हणू. मग वरील रचनेतून आपल्याला अशी
संरचना~$\Struct{M}$
मिळेल की ती~$\Gamma^*$ मधील सर्व वाक्ये पूर्ण करते, हे विगमनाने सिद्ध
करता येईल. $\Gamma \subseteq \Gamma^*$ असल्यामुळे ती~$\Gamma$ मधीलही
सर्व वाक्ये पूर्ण करेल. प्रत्यक्षात (a) आणि~(b) यांची हमी देणे पुरेसे
ठरते. ज्या वाक्यसंचासाठी (b) खरे असते त्याला \emph{संपूर्ण} म्हणतात.
त्यामुळे सुसंगत संच~$\Gamma$ चा विस्तार करून सुसंगत आणि संपूर्ण संच
$\Gamma^*$ मिळवणे हे आपले उद्दिष्ट असेल.


या योजनेत एक अडचण आहे: $\lexists[x][\varphi(x)] \in \Gamma$ असेल, तर या
प्रक्रियेत एखादे स्थिरांक~$c$ निवडून $\varphi(c)$ ची भर घालता यावी अशी
आपली अपेक्षा आहे. पण ते नेहमी करता येईल हे कसे कळणार? कदाचित आपल्या
भाषेत मोजकीच स्थिरांक असतील आणि त्यांपैकी प्रत्येकासाठी
$\lnot \varphi(c) \in \Gamma$ असेल. आपण $\varphi(c)$ चीही भर घालू शकत नाही,
कारण त्यामुळे संच विसंगत होईल; आणि $\Struct{M}$ ने $\varphi(c)$ सत्य करावे की
$\lnot \varphi(c)$ हे आपल्याला कळणार नाही. शिवाय, $\Gamma$ मध्ये अशा भाषेतील
वाक्येच असू शकतात जिच्यात एकही स्थिर चिन्ह नाही (उदा., संचसिद्धान्ताची
भाषा).

या समस्येचा उपाय म्हणजे सुरुवातीलाच अनंत स्थिरे आणि त्यांना संख्यापकांशी
योग्य प्रकारे जोडणारी वाक्ये यांची भर घालणे. (अर्थात, त्यामुळे विसंगती
निर्माण होऊ शकत नाही, हे आपल्याला तपासावे लागेल.)

आण्विक वाक्यांमध्ये फक्त स्थिरांक असतील, तर आपली मूळ रचना चांगली
चालते. पण भाषेत फलने सुद्धा असू शकतात. अशा वेळी सर्व काही योग्य
रीतीने चालण्यासाठी या फलने ना नेमून द्यायची~$\Nat$ वरील योग्य
फलने शोधणे अवघड ठरू शकते. म्हणून आणखी एक युक्ती वापरू: $\Obj c_i$ चा
अर्थ लावण्यासाठी $i$ वापरण्याऐवजी स्थिरांकाचा संचच प्रांत म्हणून
घेऊ. मग $\Struct M$ प्रत्येक स्थिरांक ला तेच स्थिर नेमून देऊ शकते:
$\Assign{\Obj c_i}{M} = \Obj c_i$. पण पूर्ण मार्गच का घेऊ नये:
$\Domain{M}$ हा भाषेतील सर्व \emph{पदांचा} संच असू द्या!{} असे केल्यास
फलने ना फलनांचे स्पष्ट निर्देशन मिळते (या फलनांची प्रचले पदे असतात
आणि त्यांची मूल्येही पदे असतात): फलन~$\Obj f^n_i$ ला आपण असे
फलन नेमून देतो की $n$ पदे $t_1$, \dots,~$t_n$ आदान म्हणून दिल्यावर ते
$\Obj f^n_i(t_1, \dots, t_n)$ हे पद मूल्य म्हणून देते.

कोड्याचा शेवटचा भाग म्हणजे~$\eq$ चे काय करायचे. विधेय~$\eq$ चा
अर्थ स्थिर आहे: $\Sat{M}{\eq[t][t']}$ तेव्हा आणि केवळ तेव्हाच, जेव्हा
$\Value{t}{M} = \Value{t'}{M}$. आता पद~$t$ चे मूल्य हे $t$ स्वतःच
असेल अशी मांडणी केली, तर $t$ आणि $t'$ हे अगदी एकच पद असल्याखेरीज ही
संरचना $\eq[t][t']$ या रूपातील \emph{एकही} वाक्य सत्य करणार नाही.
आणि ही अर्थातच समस्या आहे; कारण फलने असलेल्या भाषेतील जवळजवळ
प्रत्येक रोचक उपपत्तीमध्ये $t$ व $t'$ ही वेगळी पदे असूनही $\eq[t][t']$
या रूपाची वाक्ये प्रमेये असतील (उदा., अंकगणिताच्या उपपत्तींमध्ये
$\eq[(\Obj 0+ \Obj 0)][\Obj 0]$). ही समस्या सोडवण्यासाठी आपण
$\Struct M$ चा प्रांत बदलतो: $\Domain{M}$ मधील वस्तू म्हणून पदे वापरण्याऐवजी
पदांचे संच वापरतो; आणि प्रत्येक संचात~$\Gamma$ मधील वाक्यांनी समान असणे
आवश्यक ठरवलेली सर्व पदे असतात. उदाहरणार्थ, $\Gamma$ ही अंकगणिताची उपपत्ती
असेल, तर अशा संचांपैकी एकात $\Obj 0$, $(\Obj 0 + \Obj 0)$, $(\Obj
0 \times \Obj 0)$, इत्यादी असतील. हाच संच आपण $\Obj 0$ ला नेमून देऊ;
आणि तो त्यातील सर्व पदांचे मूल्यही आहे, उदा., $(\Obj 0 + \Obj 0)$ चेही,
हे निष्पन्न होईल. म्हणून सुधारित संरचनांमध्ये
$\eq[(\Obj 0+ \Obj 0)][\Obj 0]$ हे वाक्य सत्य असेल.

मग आपण पुढीलप्रमाणे काम करू. प्रथम संपूर्ण सुसंगत संचांच्या गुणधर्मांचा
अभ्यास करू. विशेषतः, संपूर्ण सुसंगत संचात $\varphi \land \psi$ तेव्हा आणि
केवळ तेव्हाच असते, जेव्हा त्यात $\varphi$ व~$\psi$ ही दोन्ही असतात; $\varphi \lor \psi$
तेव्हा आणि केवळ तेव्हाच असते, जेव्हा त्यांपैकी किमान एक असते; इत्यादी,
हे सिद्ध करू (\ref{fol:com:ccs:prop:ccs}). त्यानंतर वाक्यांच्या
  “संतृप्त” संचांची व्याख्या करून त्यांचा अभ्यास करू. प्रत्येक संख्यकीकृत
  वाक्याला तिच्या दृष्टांतांशी जोडणाऱ्या सशर्तांचा समावेश असलेला
  संच म्हणजे संतृप्त संच (\ref{fol:com:hen:defn:henkin-exp}). कोणत्याही सुसंगत
  संच~$\Gamma$ चा विस्तार करून संतृप्त संच~$\Gamma'$ नेहमी मिळवता येतो,
  हे आपण दाखवू (\ref{fol:com:hen:lem:henkin}). एखादा संच सुसंगत, संतृप्त आणि
  संपूर्ण असेल, तर त्यात $\lexists[x][\varphi(x)]$ तेव्हा
    आणि केवळ तेव्हाच असते, जेव्हा एखाद्या बंद पद~$t$ साठी त्यात $\varphi(t)$
    असते आणि
  $\lforall[x][\varphi(x)]$ तेव्हा आणि केवळ तेव्हाच असते,
    जेव्हा प्रत्येक बंद पद~$t$ साठी त्यात $\varphi(t)$ असते
  (\ref{fol:com:hen:prop:saturated-instances}), हाही गुणधर्म त्याला असतो.
त्यानंतर संतृप्त सुसंगत संच
$\Gamma'$ घेऊन त्याचा विस्तार करून
संतृप्त, सुसंगत आणि संपूर्ण संच~$\Gamma^*$
मिळवता येतो, हे दाखवू (\ref{fol:com:lin:lem:lindenbaum}). याच~$\Gamma^*$ चा
वापर आपण आपले पद-प्रतिमान~$\Struct M(\Gamma^*)$ परिभाषित करण्यासाठी करू.
पद-प्रतिमानाचा प्रांत हा बंद पदांचा संच असतो आणि त्यातील
  विधेयांचा अर्थ आण्विक वाक्ये यांच्या~$\Gamma^*$ मधील सदस्यत्वाने ठरते
(\ref{fol:com:mod:defn:termmodel}). संतृप्त, संपूर्ण सुसंगत
संचांचे गुणधर्म वापरून
$\Sat{M(\Gamma^*)}{\varphi}$ तेव्हा आणि
केवळ तेव्हाच, जेव्हा $\varphi \in \Gamma^*$, हे दाखवू
(\ref{fol:com:mod:lem:truth}); म्हणून विशेषतः
$\Sat{M(\Gamma^*)}{\Gamma}$.
शेवटी, $\Gamma$ मध्ये~$\eq$ सुद्धा असेल तर पद-प्रतिमानाची
  व्याख्या कशी करायची याचा विचार करू
  (\ref{fol:com:ide:defn:term-model-factor}) आणि ते~$\Gamma^*$ ची पूर्ति करते,
  हे दाखवू (\ref{fol:com:ide:lem:truth}).
















\section{वाक्याचे संपूर्ण सुसंगत संच}\label{fol:com:ccs:sec}

\begin{defn}[संपूर्ण संच]
\label{fol:com:ccs:def:complete-set} वाक्यांचा संच~$\Gamma$
\emph{संपूर्ण} असतो तेव्हा आणि केवळ तेव्हाच, जेव्हा कोणत्याही
वाक्य~$\varphi$ साठी $\varphi \in \Gamma$ किंवा $\lnot \varphi \in \Gamma$.
\end{defn}

\begin{explain}
संपूर्ण वाक्यसंच कोणताही प्रश्न अनुत्तरित ठेवत नाहीत. कोणत्याही
वाक्य~$\varphi$ साठी $\varphi$ सत्य आहे की असत्य हे~$\Gamma$ “सांगतो.”
संपूर्ण संचांचे महत्त्व संपूर्णता प्रमेयाच्या सिद्धतेपलीकडेही आहे.
उदाहरणार्थ, जी उपपत्ती संपूर्ण असून स्वयंसिद्धकीय रीतीने मांडता
येते, तिच्याविषयी नेहमी निर्णय घेता येतो.
\end{explain}

\begin{explain}
संपूर्णतेच्या सिद्धतेत संपूर्ण सुसंगत संच महत्त्वाचे असतात, कारण
वाक्यांचा प्रत्येक सुसंगत संच~$\Gamma$ हा एखाद्या संपूर्ण
सुसंगत संच~$\Gamma^*$ मध्ये समाविष्ट असतो, याची हमी आपण देऊ शकतो.
संपूर्ण सुसंगत संचात प्रत्येक वाक्य~$\varphi$ साठी $\varphi$ किंवा
त्याचे नकरण $\lnot \varphi$ यांपैकी एक असते, पण दोन्ही नसतात. हे विशेषतः
आण्विक वाक्ये यांच्यासाठी खरे असते.
म्हणून स्थिरांकाची योग्य भर घालून विस्तारित केलेल्या
भाषेतील संपूर्ण सुसंगत संचापासून आपण अशी
संरचना रचू शकतो की तिच्यातील
विधेयांचा अर्थ
हा~$\Gamma^*$ मध्ये कोणती आण्विक
वाक्ये आहेत त्यानुसार परिभाषित
केला जातो. त्यानंतर ही संरचना
$\Gamma^*$ मधील सर्व वाक्ये (आणि म्हणून~$\Gamma$ मधीलही सर्व)
सत्य करते, हे दाखवता येते. या शेवटच्या तथ्याच्या सिद्धतेसाठी
$\lnot \varphi \in \Gamma^*$ तेव्हा आणि केवळ तेव्हाच, जेव्हा $\varphi \notin
\Gamma^*$; $(\varphi \lor \psi) \in \Gamma^*$ तेव्हा आणि केवळ तेव्हाच, जेव्हा
$\varphi \in \Gamma^*$ किंवा $\psi \in \Gamma^*$; इत्यादी तथ्ये आवश्यक असतात.
\end{explain}

पुढे आपण $\Proves$ ची परावर्तिता, एकस्वनिकता आणि संक्रमकता हे गुणधर्म
अनेकदा अध्याहृतपणे वापरू (पाहा

  \ref{fol:seq:ptn:sec}, \ref{fol:ntd:ptn:sec}, \ref{fol:axd:ptn:sec}, \ref{fol:tab:ptn:sec}).

\begin{prop}
\label{fol:com:ccs:prop:ccs}
$\Gamma$ हा संपूर्ण आणि सुसंगत आहे असे समजा. मग:
\begin{enumerate}
\item \label{fol:com:ccs:prop:ccs-prov-in} जर $\Gamma \Proves \varphi$, तर $\varphi \in
  \Gamma$.

\item \label{fol:com:ccs:prop:ccs-and} $\varphi \land \psi \in \Gamma$
  तेव्हा आणि केवळ तेव्हाच, जेव्हा $\varphi \in \Gamma$ आणि $\psi \in \Gamma$
  ही दोन्ही खरे आहेत.

\item \label{fol:com:ccs:prop:ccs-or} $\varphi \lor \psi \in \Gamma$ तेव्हा
  आणि केवळ तेव्हाच, जेव्हा $\varphi \in \Gamma$ किंवा $\psi \in \Gamma$.

\item \label{fol:com:ccs:prop:ccs-if} $\varphi \lif \psi \in \Gamma$ तेव्हा
  आणि केवळ तेव्हाच, जेव्हा $\varphi \notin \Gamma$ किंवा $\psi \in \Gamma$.
\end{enumerate}
\end{prop}

\begin{proof}
पुढील सर्व प्रकरणांत $\Gamma$ हा संपूर्ण आणि सुसंगत आहे असे समजू.
\begin{enumerate}
\item जर $\Gamma \Proves \varphi$, तर $\varphi \in \Gamma$.

$\Gamma \Proves \varphi$ असे समजा. उलटपक्षी $\varphi \notin \Gamma$ असे गृहीत
धरा. $\Gamma$ हा संपूर्ण असल्यामुळे $\lnot \varphi \in \Gamma$. पुढील
संदर्भांनुसार

  \ref{fol:seq:prv:prop:explicit-inc}, \ref{fol:ntd:prv:prop:explicit-inc}, \ref{fol:axd:prv:prop:explicit-inc}, \ref{fol:tab:prv:prop:explicit-inc},
$\Gamma$ विसंगत आहे. हे $\Gamma$ सुसंगत असल्याच्या गृहीतकाशी विसंगत आहे.
म्हणून $\varphi \notin \Gamma$ असणे शक्य नाही; त्यामुळे $\varphi \in \Gamma$.

\item

$\varphi \land \psi \in \Gamma$ तेव्हा आणि केवळ तेव्हाच, जेव्हा $\varphi \in
\Gamma$ आणि $\psi \in \Gamma$ ही दोन्ही खरे आहेत:

पुढे जाणाऱ्या दिशेसाठी $\varphi \land \psi \in \Gamma$ असे समजा. मग
पुढील संदर्भांतील

  \ref{fol:seq:ppr:prop:provability-land}, \ref{fol:ntd:ppr:prop:provability-land}, \ref{fol:axd:ppr:prop:provability-land}, \ref{fol:tab:ppr:prop:provability-land}, बाब~(1) नुसार
$\Gamma \Proves \varphi$ आणि $\Gamma \Proves \psi$. \ref{fol:com:ccs:prop:ccs-prov-in}
नुसार $\varphi \in \Gamma$ आणि $\psi \in \Gamma$, अपेक्षेप्रमाणे.

उलट दिशेसाठी $\varphi \in \Gamma$ आणि $\psi \in \Gamma$ असे समजा. पुढील
संदर्भांतील

  \ref{fol:seq:ppr:prop:provability-land}, \ref{fol:ntd:ppr:prop:provability-land}, \ref{fol:axd:ppr:prop:provability-land}, \ref{fol:tab:ppr:prop:provability-land}, बाब~(2) नुसार
$\Gamma \Proves \varphi \land \psi$. \ref{fol:com:ccs:prop:ccs-prov-in} नुसार $\varphi \land
\psi \in \Gamma$.

\item

प्रथम $\varphi \lor \psi \in \Gamma$ असेल, तर $\varphi \in \Gamma$ किंवा $\psi \in
\Gamma$, हे दाखवू. $\varphi \lor \psi \in \Gamma$, पण $\varphi \notin \Gamma$ आणि
$\psi \notin \Gamma$ असे समजा. $\Gamma$ हा संपूर्ण असल्यामुळे
$\lnot \varphi \in \Gamma$ आणि $\lnot \psi \in \Gamma$. पुढील संदर्भांतील

  \ref{fol:seq:ppr:prop:provability-lor}, \ref{fol:ntd:ppr:prop:provability-lor}, \ref{fol:axd:ppr:prop:provability-lor}, \ref{fol:tab:ppr:prop:provability-lor},
बाब~(1) नुसार $\Gamma$ विसंगत आहे; हा विरोधाभास आहे. म्हणून $\varphi \in
\Gamma$ किंवा $\psi \in \Gamma$.

उलट दिशेसाठी $\varphi \in \Gamma$ किंवा $\psi \in \Gamma$ असे समजा. पुढील
संदर्भांतील

  \ref{fol:seq:ppr:prop:provability-lor}, \ref{fol:ntd:ppr:prop:provability-lor}, \ref{fol:axd:ppr:prop:provability-lor}, \ref{fol:tab:ppr:prop:provability-lor}, बाब~(2) नुसार
$\Gamma \Proves \varphi \lor \psi$. \ref{fol:com:ccs:prop:ccs-prov-in} नुसार $\varphi \lor
\psi \in \Gamma$, अपेक्षेप्रमाणे.

\item

पुढे जाणाऱ्या दिशेसाठी $\varphi \lif \psi \in \Gamma$ असे समजा आणि उलटपक्षी
$\varphi \in \Gamma$ व $\psi \notin \Gamma$ असे गृहीत धरा. या गृहीतकांनुसार
$\varphi \lif \psi \in \Gamma$ आणि $\varphi \in \Gamma$. पुढील संदर्भांतील

  \ref{fol:seq:ppr:prop:provability-lif}, \ref{fol:ntd:ppr:prop:provability-lif}, \ref{fol:axd:ppr:prop:provability-lif}, \ref{fol:tab:ppr:prop:provability-lif}, बाब~(1) नुसार
$\Gamma \Proves \psi$. पण मग \ref{fol:com:ccs:prop:ccs-prov-in} नुसार $\psi \in
\Gamma$; हे $\psi \notin \Gamma$ या गृहीतकाशी विसंगत आहे.

उलट दिशेसाठी प्रथम $\varphi \notin \Gamma$ हे प्रकरण पाहू. $\Gamma$ हा
संपूर्ण असल्यामुळे $\lnot \varphi \in \Gamma$. पुढील संदर्भांतील

  \ref{fol:seq:ppr:prop:provability-lif}, \ref{fol:ntd:ppr:prop:provability-lif}, \ref{fol:axd:ppr:prop:provability-lif}, \ref{fol:tab:ppr:prop:provability-lif}, बाब~(2) नुसार
$\Gamma \Proves \varphi \lif \psi$. पुन्हा \ref{fol:com:ccs:prop:ccs-prov-in} नुसार
$\varphi \lif \psi \in \Gamma$, अपेक्षेप्रमाणे.

आता $\psi \in \Gamma$ हे प्रकरण पाहू. पुढील संदर्भांतील

  \ref{fol:seq:ppr:prop:provability-lif}, \ref{fol:ntd:ppr:prop:provability-lif}, \ref{fol:tab:ppr:prop:provability-lif}, \ref{fol:axd:ppr:prop:provability-lif},
पुन्हा बाब~(2) नुसार $\Gamma \Proves \varphi \lif \psi$.
\ref{fol:com:ccs:prop:ccs-prov-in} नुसार $\varphi \lif \psi \in \Gamma$.
\end{enumerate}
\end{proof}


\begin{prob}
\ref{fol:com:ccs:prop:ccs} ची सिद्धता पूर्ण करा.
\end{prob}















\section{हेंकिन विस्तार}\label{fol:com:hen:sec}

\begin{explain}
संपूर्णता प्रमेय सिद्ध करण्यातील एक आव्हान असे आहे की संपूर्ण सुसंगत
संच~$\Gamma$ पासून आपण रचलेल्या प्रतिमानाने~$\Gamma$ मधील सर्व संख्यकीकृत
सूत्रे सत्य केली पाहिजेत. याची हमी देण्यासाठी लिऑन हेंकिन यांची एक
युक्ती वापरतो. थोडक्यात, भाषेत अनंत स्थिरांकाची भर घालून तिचा विस्तार
करणे आणि एक मुक्त चल असलेल्या प्रत्येक सूत्र~$\varphi(x)$ साठी
पुढील रूपाचे सूत्र जोडणे, ही ती युक्ती:
$\lexists[x][\varphi(x)] \lif \varphi(c)$,
येथे $c$ हे नवीन स्थिरांक पैकी एक आहे. $\Gamma$ पूर्ण करणारी
संरचना आपण रचू तेव्हा, यामुळे प्रत्येक
सत्य अस्तित्ववाची वाक्याचा साक्षी
नवीन स्थिरांमध्ये असेल, याची हमी मिळेल.
\end{explain}

\begin{prop}
\label{fol:com:hen:prop:lang-exp}
$\Gamma$ हा $\Lang L$ मध्ये सुसंगत असेल आणि $\Lang L$ मध्ये नवीन
स्थिरांकाचा गणनीय अनंत संच $\Obj d_0$, $\Obj d_1$, \dots
मिळवून $\Lang L'$ तयार केली असेल, तर $\Gamma$ हा~$\Lang L'$ मध्ये
सुसंगत असतो.
\end{prop}

\begin{defn}[संतृप्त संच]
भाषा $\Lang {L}$ मधील सूत्रांचा संच $\Gamma$
\emph{संतृप्त} असतो तेव्हा आणि केवळ तेव्हाच, जेव्हा एक मुक्त
चल~$x$ असलेल्या प्रत्येक सूत्र~$\varphi(x) \in \Frm[L]$
साठी असे स्थिरांक~$c \in \Lang{L}$ असते की
$\lexists[x][\varphi(x)] \lif \varphi(c) \in \Gamma$.
\end{defn}

पुढील प्रमेयाच्या सिद्धतेत खालील व्याख्या वापरली जाईल.

\begin{defn}
\label{fol:com:hen:defn:henkin-exp}
$\Lang L'$ ही \ref{fol:com:hen:prop:lang-exp} मधीलप्रमाणे असू द्या. ज्या
$\Lang L'$ मधील सूत्रे~$\varphi_i(x_i)$ मध्ये एक चर ($x_i$) मुक्तपणे
येतो, त्या सर्वांचे $\varphi_0(x_0)$, $\varphi_1(x_1)$, \dots असे प्रगणन निश्चित
करा. आपण~$n$ वर विगमनाने वाक्ये~$\theta_n$ परिभाषित करतो.

$\Lang{L}$ मध्ये आपण भर घातलेल्या $\Obj d_i$ पैकी~$\varphi_0(x_0)$ मध्ये
न येणारे पहिले स्थिरांक म्हणजे $c_0$ असू द्या. $\theta_0$,
\dots,~$\theta_{n-1}$ आधीच परिभाषित केली आहेत असे गृहीत धरून, नवीन
स्थिरांक~$\Obj d_i$ पैकी~$\theta_0$, \dots,~$\theta_{n-1}$ यांमध्ये किंवा
$\varphi_n(x_n)$ मध्ये न येणारे पहिले स्थिर म्हणजे $c_n$ असू द्या.

आता $\theta_{n}$ हे पुढील सूत्र असू द्या:
$\lexists[x_{n}][\varphi_{n}(x_{n})] \lif
  \varphi_{n}(c_{n})$.
\end{defn}

\begin{lem}
\label{fol:com:hen:lem:henkin}
प्रत्येक सुसंगत संच~$\Gamma$ चा विस्तार करून संतृप्त सुसंगत
संच~$\Gamma'$ मिळवता येतो.
\end{lem}

\begin{proof}
भाषा~$\Lang{L}$ मधील वाक्यांचा सुसंगत संच~$\Gamma$ दिला असता,
नवीन स्थिरांकाचा गणनीय अनंत संच जोडून भाषेचा विस्तार
$\Lang{L'}$ पर्यंत करा. \ref{fol:com:hen:prop:lang-exp} नुसार $\Gamma$ या अधिक
समृद्ध भाषेतही सुसंगत असतो. तसेच $\theta_i$ ही
\ref{fol:com:hen:defn:henkin-exp} मधीलप्रमाणे असू द्या. पुढील व्याख्या करा:
\begin{align*}
\Gamma_0 & = \Gamma \\
\Gamma_{n+1} & = \Gamma_n \cup \{\theta_n \}
\end{align*}
म्हणजे $\Gamma_{n+1} = \Gamma \cup \{ \theta_0, \dots, \theta_n \}$; आणि
$\Gamma' = \bigcup_{n} \Gamma_n$ असू द्या. $\Gamma'$ स्पष्टपणे संतृप्त आहे.

$\Gamma'$ विसंगत असेल, तर कोणत्यातरी $n$ साठी $\Gamma_n$ विसंगत असेल
(सराव: का ते स्पष्ट करा). म्हणून $\Gamma'$ सुसंगत आहे हे दाखवण्यासाठी
प्रत्येक संच~$\Gamma_n$ सुसंगत आहे हे~$n$ वर विगमनाने दाखवणे पुरेसे आहे.

विगमनाचा आधार म्हणजे $\Gamma_0 = \Gamma$ सुसंगत आहे हे विधानच; ते
प्रमेयाचे गृहीतक आहे. विगमन-पायरीसाठी $\Gamma_{n}$ सुसंगत आहे, पण
$\Gamma_{n+1} = \Gamma_n \cup \{\theta_n\}$ विसंगत आहे असे समजा. लक्षात
घ्या की $\theta_n$ हे
$\lexists[x_{n}][\varphi_{n}(x_n)] \lif \varphi_{n}(c_{n})$
आहे; येथे $\varphi_n(x_n)$ हे फक्त~$x_n$ चर मुक्त असलेले $\Lang{L'}$ चे
सूत्र आहे. $c_n$ निवडण्याच्या पद्धतीनुसार
(\ref{fol:com:hen:defn:henkin-exp} पाहा), $c_n$ हे~$\varphi_n(x_n)$ किंवा~$\Gamma_n$
मध्ये येत नाही.

$\Gamma_n \cup \{\theta_n\}$ विसंगत असेल, तर $\Gamma_n
\Proves \lnot \theta_n$; म्हणून पुढील दोन्ही विधाने खरी आहेत:

\[\Gamma_n \Proves \lexists[x_n][\varphi_n(x_n)]
\qquad
\Gamma_n \Proves \lnot \varphi_n(c_n)
\]
$c_n$ हे $\Gamma_n$ किंवा~$\varphi_n(x_n)$ मध्ये येत नसल्यामुळे
\ref{fol:axd:qpr:thm:strong-generalization}, \ref{fol:seq:qpr:thm:strong-generalization}, \ref{fol:ntd:qpr:thm:strong-generalization}, \ref{fol:tab:qpr:thm:strong-generalization}
लागू होते.
$\Gamma_n \Proves \lnot \varphi_n(c_n)$
यापासून अनुक्रमे
$\Gamma_n \Proves \lforall[x_n][\lnot \varphi_n(x_n)]$
मिळते. अशा रीतीने आपल्याकडे पुढील दोन्ही विधाने आहेत:
$\Gamma_n \Proves \lexists[x_n][\varphi_n(x_n)]$ आणि
$\Gamma_n \Proves \lforall[x_n][\lnot \varphi_n(x_n)]$;
म्हणून $\Gamma_n$ स्वतः विसंगत आहे.
(लक्षात घ्या की
$\lforall[x_n][\lnot \varphi_n(x_n)] \Proves
  \lnot\lexists[x_n][\varphi_n(x_n)]$.)
हा विरोधाभास आहे: $\Gamma_n$ सुसंगत असल्याचे गृहीत होते. म्हणून
$\Gamma_n \cup \{ \theta_n\}$ सुसंगत आहे.

\end{proof}

\begin{explain}
आता आपण दाखवू की संतृप्त असलेल्या \emph{संपूर्ण}, सुसंगत संचांना पुढील
गुणधर्म असतो: त्यात सर्ववाची संख्यकीकृत वाक्य
  तेव्हा आणि केवळ तेव्हाच असते, जेव्हा त्याचे सर्व दृष्टांत त्यात असतात
   आणि त्यात अस्तित्ववाची
  संख्यकीकृत वाक्य तेव्हा आणि केवळ तेव्हाच असते, जेव्हा त्याचा
  किमान एक दृष्टांत त्यात असतो. संपूर्ण, सुसंगत, संतृप्त संचापासून
आपण निर्माण करू ती संरचना त्यातील सर्व संख्यकीकृत वाक्ये सत्य करते,
हे दाखवण्यासाठी हा गुणधर्म वापरू.
\end{explain}

\begin{prop}\label{fol:com:hen:prop:saturated-instances}
$\Gamma$ हा संपूर्ण, सुसंगत आणि संतृप्त आहे असे समजा.
\begin{enumerate}
\item $\lexists[x][\varphi(x)] \in \Gamma$ तेव्हा आणि केवळ तेव्हाच,
  जेव्हा किमान एका बंद पद~$t$ साठी $\varphi(t) \in \Gamma$.
\item $\lforall[x][\varphi(x)] \in \Gamma$ तेव्हा आणि केवळ तेव्हाच,
  जेव्हा सर्व बंद पदे~$t$ यांच्यासाठी $\varphi(t) \in \Gamma$.
\end{enumerate}
\end{prop}

\begin{proof}
  \begin{enumerate}
  \item
    प्रथम $\lexists[x][\varphi(x)]
      \in \Gamma$ असे समजा. $\Gamma$ संतृप्त असल्यामुळे कोणत्यातरी
      स्थिरांक~$c$ साठी
      $(\lexists[x][\varphi(x)] \lif \varphi(c)) \in \Gamma$. पुढील संदर्भांतील
      \ref{fol:axd:ppr:prop:provability-lif}, \ref{fol:seq:ppr:prop:provability-lif}, \ref{fol:ntd:ppr:prop:provability-lif}, \ref{fol:tab:ppr:prop:provability-lif},
      बाब~(1), आणि
      \ref{fol:com:ccs:prop:ccs}\ref{fol:com:ccs:prop:ccs-prov-in} नुसार $\varphi(c)
      \in \Gamma$.

    उलट दिशेसाठी संतृप्तता आवश्यक नाही: $\varphi(t) \in \Gamma$ असे समजा.
    मग पुढील संदर्भांतील
      \ref{fol:axd:qpr:prop:provability-quantifiers}, \ref{fol:seq:qpr:prop:provability-quantifiers}, \ref{fol:ntd:qpr:prop:provability-quantifiers}, \ref{fol:tab:qpr:prop:provability-quantifiers},
      बाब~(1) नुसार $\Gamma \Proves \lexists[x][\varphi(x)]$.
    \ref{fol:com:ccs:prop:ccs}\ref{fol:com:ccs:prop:ccs-prov-in} नुसार
    $\lexists[x][\varphi(x)] \in \Gamma$.
  \item
    सर्व बंद पदे~$t$ यांच्यासाठी
      $\varphi(t) \in \Gamma$ असे समजा. विरोधासाठी
      $\lforall[x][\varphi(x)] \notin \Gamma$ असे गृहीत धरा. $\Gamma$ संपूर्ण
      असल्यामुळे $\lnot\lforall[x][\varphi(x)] \in \Gamma$. संतृप्ततेनुसार
      कोणत्यातरी स्थिरांक~$c$ साठी
      $(\lexists[x][\lnot \varphi(x)] \lif \lnot \varphi(c)) \in
        \Gamma$. $c$ हे बंद पद असल्यामुळे गृहीतकानुसार $\varphi(c) \in
      \Gamma$. पण यामुळे $\Gamma$ विसंगत होईल. (सराव: पुढील संच विसंगत
      असल्याचे दाखवणारी निष्पत्ती द्या:
      \[      \lnot \lforall[x][\varphi(x)],
      \lexists[x][\lnot \varphi(x)] \lif \lnot
        \varphi(c), \varphi(c)
      \]
      .)

      उलट दिशेसाठी संतृप्तता आवश्यक नाही:
      $\lforall[x][\varphi(x)] \in \Gamma$ असे समजा. मग पुढील संदर्भांतील
      \ref{fol:axd:qpr:prop:provability-quantifiers}, \ref{fol:seq:qpr:prop:provability-quantifiers}, \ref{fol:ntd:qpr:prop:provability-quantifiers}, \ref{fol:tab:qpr:prop:provability-quantifiers},
      बाब~(2) नुसार $\Gamma \Proves \varphi(t)$. \ref{fol:com:ccs:prop:ccs}
      नुसार $\varphi(t) \in \Gamma$.
  \end{enumerate}
\end{proof}












\section{लिंडेनबाउमचे पूर्वप्रमेय}\label{fol:com:lin:sec}

\begin{explain}
आता आपण असे पूर्वप्रमेय सिद्ध करू की वाक्यांचा कोणताही सुसंगत
संच केवळ सुसंगतच नव्हे, तर संपूर्ण असलेल्या एखाद्या वाक्यसंचात
समाविष्ट असतो. या सिद्धतेत प्रत्येक वेळी एक वाक्य जोडले जाते
आणि प्रत्येक पायरीवर संच सुसंगत राहील याची हमी दिली जाते. हे अशा रीतीने
करतो की प्रत्येक $\varphi$ साठी कोणत्यातरी टप्प्यावर $\varphi$ किंवा $\lnot
\varphi$ यांपैकी एक जोडले जाईल. मग त्या रचनेतील सर्व टप्प्यांचा संघ $\varphi$
किंवा त्याचे नकरण~$\lnot \varphi$ यांपैकी एक समाविष्ट करतो आणि म्हणून तो
संपूर्ण असतो. प्रत्येक टप्प्यावर विसंगती येणार नाही याची खात्री केल्यामुळे
तो सुसंगतही असतो.
\end{explain}

\begin{lem}[लिंडेनबाउमचे पूर्वप्रमेय]
\label{fol:com:lin:lem:lindenbaum} भाषा~$\Lang{L}$ मधील प्रत्येक सुसंगत
संच~$\Gamma$ चा विस्तार करून संपूर्ण आणि सुसंगत असा
संच~$\Gamma^*$ मिळवता येतो.
\end{lem}

\begin{proof}
$\Gamma$ सुसंगत असू द्या. $\Lang L$ मधील सर्व वाक्यांचे
$\varphi_0$, $\varphi_1$, \dots{} असे प्रगणन असू द्या. $\Gamma_0 = \Gamma$
अशी व्याख्या करा, आणि
\[\Gamma_{n+1} =
\begin{cases}
\Gamma_n \cup \{ \varphi_n \} & \textrm{जर $\Gamma_n \cup \{\varphi_n\}$
  सुसंगत असेल;} \\
\Gamma_n \cup \{ \lnot \varphi_n \} & \textrm{अन्यथा.}
\end{cases}
\]
$\Gamma^* = \bigcup_{n \geq 0} \Gamma_n$ असू द्या.

प्रत्येक $\Gamma_n$ सुसंगत आहे: व्याख्येनुसार $\Gamma_0$ सुसंगत आहे.
$\Gamma_{n+1} = \Gamma_n \cup \{\varphi_n\}$ असेल, तर याचे कारण हा
दुसरा संच सुसंगत असणे हेच आहे. तो सुसंगत नसेल, तर $\Gamma_{n+1} =
\Gamma_n \cup \{\lnot \varphi_n\}$. $\Gamma_n \cup \{\lnot \varphi_n\}$
सुसंगत आहे हे आपल्याला पडताळावे लागेल. तो सुसंगत नाही असे समजा. मग
$\Gamma_n \cup \{\varphi_n\}$ आणि $\Gamma_n \cup \{\lnot \varphi_n\}$ हे
\emph{दोन्ही} विसंगत आहेत. याचा अर्थ

  \ref{fol:axd:prv:prop:provability-exhaustive}, \ref{fol:seq:prv:prop:provability-exhaustive}, \ref{fol:ntd:prv:prop:provability-exhaustive}, \ref{fol:tab:prv:prop:provability-exhaustive}
नुसार $\Gamma_n$ विसंगत असेल; पण हे विगमन-गृहीतकाच्या विरुद्ध आहे.

प्रत्येक~$n$ आणि प्रत्येक $i < n$ साठी $\Gamma_i \subseteq \Gamma_n$.
हे~$n$ वरील सोप्या विगमनाने मिळते. $n=0$ असेल, तर $i < 0$ असा
कोणताही निर्देशांक नसतो; म्हणून विधान आपोआप लागू होते. विगमन-पायरीसाठी ते~$n$
साठी खरे आहे असे समजा. $i < n+1$ असेल, तर $\Gamma_i \subseteq
\Gamma_{n+1}$ हे दाखवू. रचनेनुसार $\Gamma_{n+1} = \Gamma_n \cup
\{\varphi_n\}$ किंवा $= \Gamma_n \cup \{\lnot \varphi_n\}$. म्हणून $\Gamma_n
\subseteq \Gamma_{n+1}$. $i < n+1$ असेल, तर विगमन-गृहीतकानुसार
$\Gamma_i \subseteq \Gamma_n$ ($i < n$ असल्यास), किंवा $\Gamma_n
\subseteq \Gamma_n$ हे क्षुल्लक तथ्य लागू होते ($i = n$ असल्यास).
म्हणून~$\subseteq$ च्या संक्रमणशीलतेने $\Gamma_i \subseteq
\Gamma_{n+1}$ मिळते.

यावरून $\Gamma^*$ सुसंगत आहे हे मिळते. ते कसे ते पाहू: $\Gamma'
\subseteq \Gamma^*$ हा सांत संच असू द्या. प्रत्येक $\psi \in \Gamma'$
हे कोणत्यातरी~$i$ साठी~$\Gamma_i$ मध्येही असते. या निर्देशांकांतील
सर्वांत मोठा निर्देशांक $n$ असू द्या. $i \le n$ असेल, तर $\Gamma_i
\subseteq \Gamma_n$ असल्यामुळे प्रत्येक $\psi \in \Gamma'$ हे
$\in \Gamma_n$ देखील आहे, म्हणजे $\Gamma' \subseteq \Gamma_n$, आणि
$\Gamma_n$ सुसंगत आहे. म्हणून प्रत्येक सांत उपसंच $\Gamma' \subseteq
\Gamma^*$ सुसंगत आहे.

  \ref{fol:axd:ptn:prop:proves-compact}, \ref{fol:seq:ptn:prop:proves-compact}, \ref{fol:ntd:ptn:prop:proves-compact}, \ref{fol:tab:ptn:prop:proves-compact} नुसार $\Gamma^*$
सुसंगत आहे.

$\Frm[L]$ मधील प्रत्येक वाक्य हा~$\Gamma^*$ ची व्याख्या
करण्यासाठी वापरलेल्या यादीत येतो. $\varphi_n \notin \Gamma^*$ असेल, तर
$\Gamma_n \cup \{\varphi_n\}$ विसंगत होते हेच त्याचे कारण आहे. पण मग
$\lnot \varphi_n \in \Gamma^*$; म्हणून $\Gamma^*$ हा संपूर्ण आहे.
\end{proof}












\section{प्रतिमानाची रचना}\label{fol:com:mod:sec}

\begin{explain}
सध्या आपण~$\eq$ चा विचार करत नाही; म्हणजे~$\eq$ नसलेला
  वाक्यांचा सुसंगत संच~$\Gamma$ पूर्ततायोग्य आहे एवढेच आपल्याला
  दाखवायचे आहे. प्रथम~$\Gamma$ चा विस्तार करून सुसंगत, संपूर्ण
  आणि संतृप्त असा संच~$\Gamma^*$ मिळवतो. या बाबतीत
  प्रतिमान~$\Struct{M(\Gamma^*)}$ ची व्याख्या सोपी आहे: आपण~$\Lang{L'}$
  मधील बंद पदांचा संच प्रांत म्हणून घेतो. प्रत्येक स्थिरांक ला तेच
  स्थिर नेमतो आणि, अधिक व्यापकपणे, प्रत्येक बंद पद~$t$ साठी
  $\Value{t}{M(\Gamma^*)} = t$ याची खात्री करतो. विधेये ना
  असे विस्तार नेमतो की एखादे आण्विक वाक्य हे
  $\Struct{M(\Gamma^*)}$ मध्ये तेव्हा आणि केवळ तेव्हाच सत्य असते, जेव्हा
  ते~$\Gamma^*$ मध्ये असते. त्यामुळे~$\Gamma^*$ मधील सर्व आण्विक
  वाक्ये हे~$\Struct{M(\Gamma^*)}$ मध्ये सत्य होतील, हे उघड आहे.
  आपण सुरुवात केलेला~$\Gamma^*$ सुसंगत, संपूर्ण आणि संतृप्त असेल, तर
  उरलेली वाक्येही सत्य असतात.
\end{explain}


\begin{defn}[पद-प्रतिमान]
\label{fol:com:mod:defn:termmodel}
$\Gamma^*$ हा भाषा~$\Lang L$ मधील वाक्यांचा संपूर्ण,
सुसंगत आणि संतृप्त संच असू द्या. $\Gamma^*$ चे
\emph{पद-प्रतिमान}~$\Struct M(\Gamma^*)$ ही पुढीलप्रमाणे परिभाषित केलेली
संरचना आहे:
\begin{enumerate}
\item प्रांत~$\Domain{M(\Gamma^*)}$ हा~$\Lang L$ मधील सर्व बंद
  पदांचा संच आहे.
\item स्थिरांक~$c$ चा निर्वचनार्थ स्वतः~$c$ आहे:
  $\Assign{c}{M(\Gamma^*)} = c$.
\item फलन~$f$ ला असे फलन नेमले जाते की बंद पदे $t_1$, \dots,
  $t_n$ ही वितर्के दिली असता त्याचे मूल्य बंद पद $f(t_1, \dots, t_n)$
  असते:
\[\Assign{f}{M(\Gamma^*)}(t_1, \dots, t_n) = f(t_1,\dots, t_n)
\]
\item $R$ हे $n$-स्थानी विधेय असेल, तर
  \[  \tuple{t_1, \dots,
  t_n} \in \Assign{R}{M(\Gamma^*)} \text{ तेव्हा आणि केवळ तेव्हाच } \Atom{R}{t_1, \dots,
    t_n} \in \Gamma^*.
  \]
\end{enumerate}
\end{defn}



आपल्याकडे खरोखरच $\Value{t}{M(\Gamma^*)} = t$ आहे, हे आता पडताळू.

\begin{lem}
\label{fol:com:mod:lem:val-in-termmodel} \ref{fol:com:mod:defn:termmodel} मधील
$\Struct M(\Gamma^*)$ हे पद-प्रतिमान असू द्या; मग
$\Value{t}{M(\Gamma^*)} = t$.
\end{lem}
\begin{proof}
 ही सिद्धता~$t$ वरील विगमनाने मिळते. $t$ हे स्थिरांक असण्याचा
 आधार-प्रसंग पद-प्रतिमानाच्या व्याख्येवरून थेट मिळतो. विगमन-पायरीसाठी
 $t_1, \ldots, t_n$ ही अशी बंद पदे आहेत की
 $\Value{t_i}{M(\Gamma^*)} = t_i$, आणि $f$ हे $n$-स्थानी
 फलन आहे, असे समजा. मग
\begin{align*}
\Value{f(t_1,\ldots,t_n)}{M(\Gamma^*)} &= \Assign{f}{M(\Gamma^*)}(\Value{t_1}
 {M(\Gamma^*)},
\ldots, \Value{t_n}{M(\Gamma^*)}) \\
&= \Assign{f}{M(\Gamma^*)}(t_1, \dots, t_n) \\
&= f(t_1,\dots, t_n),
\end{align*}
आणि म्हणून विगमनाने हे प्रत्येक बंद पद~$t$ साठी लागू होते.
\end{proof}



\begin{explain}
  एखादी संरचना~$\Struct{M}$ अस्तित्ववाची संख्यकीकृत
    वाक्य~$\lexists[x][\varphi(x)]$ सत्य करू शकते, पण ती सत्य करते
    असा कोणताही दृष्टांत~$\varphi(t)$ नसेल.
  एखादी संरचना~$\Struct{M}$ सर्ववाची संख्यकीकृत
    वाक्य~$\lforall[x][\varphi(x)]$ चे सर्व दृष्टांत~$\varphi(t)$ सत्य करू
    शकते, पण~$\lforall[x][\varphi(x)]$ सत्य करणार नाही. याचे कारण सामान्यतः
  प्रांत~$\Domain{M}$ मधील प्रत्येक घटक हे एखाद्या बंद पदाचे
  मूल्य नसते ($\Struct{M}$ बंद पदांनी आच्छादित नसेल). म्हणूनच
  पूर्ति-संबंधाची व्याख्या चर-मूल्यांकनांच्या साहाय्याने केली जाते. पण आपल्या
  पद-प्रतिमान~$\Struct{M(\Gamma^*)}$ साठी त्याची आवश्यकता नसती---कारण
  ते बंद पदांनी आच्छादित आहे. पुढील निष्पत्तीचा आशय हाच आहे.
\end{explain}

\begin{prop}
\label{fol:com:mod:prop:quant-termmodel}
\ref{fol:com:mod:defn:termmodel} मधील $\Struct M(\Gamma^*)$ हे पद-प्रतिमान असू द्या.
\begin{enumerate}
\item $\Sat{M(\Gamma^*)}{\lexists[x][\varphi(x)]}$ तेव्हा आणि केवळ तेव्हाच,
  जेव्हा किमान एका बंद पद~$t$ साठी $\Sat{M(\Gamma^*)}{\varphi(t)}$.
\item $\Sat{M(\Gamma^*)}{\lforall[x][\varphi(x)]}$ तेव्हा आणि केवळ तेव्हाच,
  जेव्हा सर्व बंद पदे~$t$ यांच्यासाठी $\Sat{M(\Gamma^*)}{\varphi(t)}$.
\end{enumerate}
\end{prop}

\begin{proof}
\begin{enumerate}
\item
  \ref{fol:syn:ass:prop:sat-quant} नुसार
    $\Sat{M(\Gamma^*)}{\lexists[x][\varphi(x)]}$ तेव्हा आणि केवळ तेव्हाच,
    जेव्हा किमान एका चर-मूल्यांकन~$s$ साठी
    $\Sat{M(\Gamma^*)}{\varphi(x)}[s]$. $\Domain{M(\Gamma^*)}$ मध्ये
    $\Lang{L}$ मधील बंद पदे असल्यामुळे, हे तेव्हा आणि केवळ तेव्हाच लागू
    होते, जेव्हा असे किमान एक बंद पद~$t$ असते की $s(x) = t$ आणि
    $\Sat{M(\Gamma^*)}{\varphi(x)}[s]$. \ref{fol:syn:ext:prop:ext-formulas}
    नुसार, $s(x) = t$ असेल तेथे $\Sat{M(\Gamma^*)}{\varphi(x)}[s]$ तेव्हा आणि
    केवळ तेव्हाच, जेव्हा $\Sat{M(\Gamma^*)}{\varphi(t)}[s]$. तसेच
    \ref{fol:syn:ass:prop:sentence-sat-true} नुसार,
    $\varphi(t)$ हे वाक्य असल्यामुळे $\Sat{M(\Gamma^*)}{\varphi(t)}[s]$ तेव्हा आणि
    केवळ तेव्हाच, जेव्हा $\Sat{M(\Gamma^*)}{\varphi(t)}$.
\item
  \ref{fol:syn:ass:prop:sat-quant} नुसार
    $\Sat{M(\Gamma^*)}{\lforall[x][\varphi(x)]}$ तेव्हा आणि केवळ तेव्हाच,
    जेव्हा प्रत्येक चर-मूल्यांकन $s$ साठी
    $\Sat{M(\Gamma^*)}{\varphi(x)}[s]$. लक्षात घ्या की
    $\Domain{M(\Gamma^*)}$ मध्ये~$\Lang{L}$ मधील बंद पदे असतात. त्यामुळे
    प्रत्येक बंद पद~$t$ साठी $s(x) = t$ असे चर-मूल्यांकन मिळते, आणि
    कोणत्याही चर-मूल्यांकनासाठी $s(x)$ हे एखादे बंद पद~$t$ असते.
    \ref{fol:syn:ext:prop:ext-formulas} नुसार, $s(x) = t$ असेल तेथे
    $\Sat{M(\Gamma^*)}{\varphi(x)}[s]$ तेव्हा आणि केवळ तेव्हाच, जेव्हा
    $\Sat{M(\Gamma^*)}{\varphi(t)}[s]$. तसेच
    \ref{fol:syn:ass:prop:sentence-sat-true} नुसार,
    $\varphi(t)$ हे वाक्य असल्यामुळे $\Sat{M(\Gamma^*)}{\varphi(t)}[s]$ तेव्हा आणि
    केवळ तेव्हाच, जेव्हा $\Sat{M(\Gamma^*)}{\varphi(t)}$.
\end{enumerate}
\end{proof}




\begin{lem}[सत्यता पूर्वप्रमेय]
\label{fol:com:mod:lem:truth} $\varphi$ मध्ये~$\eq$ येत नाही असे समजा. मग
$\Sat{M(\Gamma^*)}{\varphi}$ तेव्हा आणि केवळ तेव्हाच, जेव्हा $\varphi \in \Gamma^*$.
\end{lem}

\begin{proof}
दोन्ही दिशा एकाच वेळी, $\varphi$ वरील विगमनाने सिद्ध करू.
\begin{enumerate}
\item \indcase{\varphi}{\lfalse}
  {$\Sat/{M(\Gamma^*)}{\lfalse}$
    हे पूर्तिच्या व्याख्येनुसार मिळते. दुसरीकडे, $\Gamma^*$ सुसंगत
    असल्यामुळे $\lfalse \notin \Gamma^*$.}

\item \indcase{\varphi}{\ltrue}
  {$\Sat{M(\Gamma^*)}{\ltrue}$
    हे पूर्तिच्या व्याख्येनुसार मिळते. दुसरीकडे, $\Gamma^*$ सुसंगत व
    संपूर्ण असल्यामुळे आणि $\Gamma^* \Proves \ltrue$ असल्यामुळे
    $\ltrue \in \Gamma^*$.}

\item \indcase{\varphi}{R(t_1, \dots, t_n)}
  {$\Sat{M(\Gamma^*)}{\Atom{R}{t_1, \dots, t_n}}$ तेव्हा आणि केवळ तेव्हाच,
    जेव्हा $\tuple{t_1, \dots, t_n} \in \Assign{R}{M(\Gamma^*)}$
    (पूर्तिच्या व्याख्येनुसार); आणि हे तेव्हा आणि केवळ तेव्हाच, जेव्हा
    $R(t_1, \dots, t_n) \in \Gamma^*$ ($\Struct M(\Gamma^*)$ च्या
    रचनेनुसार).}

\item
  \indcase{\varphi}{\lnot \psi}
    {$\Sat{M(\Gamma^*)}{\indfrm}$ तेव्हा
    आणि केवळ तेव्हाच, जेव्हा
    $\Sat/{M(\Gamma^*)}{\psi}$
    (पूर्तिच्या व्याख्येनुसार). विगमन-गृहीतकानुसार,
    $\Sat/{M(\Gamma^*)}{\psi}$ तेव्हा
    आणि केवळ तेव्हाच, जेव्हा $\psi \notin \Gamma^*$. $\Gamma^*$ सुसंगत
    आणि संपूर्ण असल्यामुळे $\psi \notin \Gamma^*$ तेव्हा आणि केवळ
    तेव्हाच, जेव्हा $\lnot \psi \in \Gamma^*$.}

\item

  \indcase{\varphi}{\psi \land
    \chi}{$\Sat{M(\Gamma^*)}{\indfrm}$
    तेव्हा आणि केवळ तेव्हाच, जेव्हा
    $\Sat{M(\Gamma^*)}{\psi}$ आणि
    $\Sat{M(\Gamma^*)}{\chi}$ हे दोन्ही
    लागू होतात (पूर्तिच्या व्याख्येनुसार); आणि हे तेव्हा आणि केवळ तेव्हाच,
    जेव्हा $\psi \in \Gamma^*$ आणि $\chi \in \Gamma^*$ हे दोन्ही लागू
    होतात (विगमन-गृहीतकानुसार). तसेच
    \ref{fol:com:ccs:prop:ccs}\ref{fol:com:ccs:prop:ccs-and} नुसार, हे तेव्हा आणि
    केवळ तेव्हाच लागू होते, जेव्हा $(\psi \land \chi) \in \Gamma^*$.}

\item

  \indcase{\varphi}{\psi \lor \chi}
    {$\Sat{M(\Gamma^*)}{\indfrm}$
    तेव्हा आणि केवळ तेव्हाच, जेव्हा
    $\Sat{M(\Gamma^*)}{\psi}$ किंवा
    $\Sat{M(\Gamma^*)}{\chi}$
    (पूर्तिच्या व्याख्येनुसार); आणि हे तेव्हा आणि केवळ तेव्हाच, जेव्हा
    $\psi \in \Gamma^*$ किंवा $\chi \in \Gamma^*$ (विगमन-गृहीतकानुसार).
    \ref{fol:com:ccs:prop:ccs}\ref{fol:com:ccs:prop:ccs-or} नुसार, हे तेव्हा आणि
    केवळ तेव्हाच लागू होते, जेव्हा $(\psi \lor \chi) \in \Gamma^*$.}

\item

  \indcase{\varphi}{\psi \lif \chi}
    {$\Sat{M(\Gamma^*)}{\indfrm}$
    तेव्हा आणि केवळ तेव्हाच, जेव्हा
    $\Sat/{M(\Gamma^*)}{\psi}$ किंवा $\Sat{M (\Gamma^*)}{\chi}$
    (पूर्तिच्या व्याख्येनुसार); आणि हे तेव्हा आणि केवळ तेव्हाच, जेव्हा
    $\psi \notin \Gamma^*$ किंवा $\chi \in \Gamma^*$ (विगमन-गृहीतकानुसार).
    \ref{fol:com:ccs:prop:ccs}\ref{fol:com:ccs:prop:ccs-if} नुसार, हे तेव्हा आणि
    केवळ तेव्हाच लागू होते, जेव्हा $(\psi \lif \chi) \in \Gamma^*$.}


\item

    \indcase{\varphi}{\lforall[x][\psi(x)]}{$\Sat{M(\Gamma^*)}{\indfrm}$ तेव्हा
      आणि केवळ तेव्हाच, जेव्हा सर्व बंद पदे~$t$ यांच्यासाठी
      $\Sat{M(\Gamma^*)}{\psi(t)}$ (\ref{fol:com:mod:prop:quant-termmodel}).
      विगमन-गृहीतकानुसार हे तेव्हा आणि केवळ तेव्हाच लागू होते, जेव्हा सर्व
      बंद पदे~$t$ यांच्यासाठी $\psi(t) \in \Gamma^*$; आणि
      \ref{fol:com:hen:prop:saturated-instances} नुसार हे पुन्हा तेव्हा आणि
      केवळ तेव्हाच लागू होते, जेव्हा $\lforall[x][\psi(x)] \in \Gamma^*$.}


\item

    \indcase{\varphi}{\lexists[x][\psi(x)]}{$\Sat{M(\Gamma^*)}{\indfrm}$ तेव्हा
      आणि केवळ तेव्हाच, जेव्हा किमान एका बंद पद~$t$ साठी
      $\Sat{M(\Gamma^*)}{\psi(t)}$ (\ref{fol:com:mod:prop:quant-termmodel}).
      विगमन-गृहीतकानुसार हे तेव्हा आणि केवळ तेव्हाच लागू होते, जेव्हा किमान
      एका बंद पद~$t$ साठी $\psi(t) \in \Gamma^*$. तसेच
      \ref{fol:com:hen:prop:saturated-instances} नुसार हे पुन्हा तेव्हा आणि
      केवळ तेव्हाच लागू होते, जेव्हा $\lexists[x][\psi(x)] \in \Gamma^*$.}


\end{enumerate}
\end{proof}
















\section{एकरूपता}\label{fol:com:ide:sec}

\begin{explain}
मागील विभागात दिलेली पद-प्रतिमानाची रचना ही~$\eq$ नसलेल्या
संच~$\Gamma$ साठी प्रथम-कोटी तर्कशास्त्राची संपूर्णता प्रस्थापित करण्यास
पुरेशी आहे. पद-प्रतिमान हे~$\eq$ नसलेले प्रत्येक $\varphi \in \Gamma^*$
पूर्त करते (आणि म्हणून सर्व $\varphi \in \Gamma$ पूर्त करते). पण~$\eq$
उपस्थित असेल, तर ही रचना चालत नाही. कारण मग~$\Gamma^*$ मध्ये
वाक्य~$\eq[t][t']$ असू शकते, पण पद-प्रतिमानात कोणत्याही पदाचे
मूल्य ते पद स्वतःच असते. म्हणून $t$ आणि $t'$ ही भिन्न पदे असतील, तर
पद-प्रतिमानातील त्यांची मूल्ये---म्हणजे अनुक्रमे $t$ आणि $t'$---भिन्न
असतात; त्यामुळे $\eq[t][t']$ असत्य असते. मात्र ``भागाकार घेणे'' म्हणून
ओळखली जाणारी रचना वापरून हे दुरुस्त करता येते.
\end{explain}

\begin{defn}
  $\Gamma^*$ हा~$\Lang L$ मधील वाक्यांचा सुसंगत आणि संपूर्ण संच
  असू द्या. $\Lang L$ मधील बंद पदांच्या संचावर आपण~$\approx$ हा संबंध
  पुढीलप्रमाणे परिभाषित करतो:
  \[  t \approx t' \text{\quad तेव्हा आणि केवळ तेव्हाच \quad} \eq[t][t'] \in \Gamma^*
  \]
\end{defn}

\begin{prop}
\label{fol:com:ide:prop:approx-equiv}
$\approx$ या संबंधाला पुढील गुणधर्म आहेत:
\begin{enumerate}
\item $\approx$ परावर्ती आहे.
\item $\approx$ सममित आहे.
\item $\approx$ संक्रमक आहे.
\item $t \approx t'$, $f$ हे फलन, आणि $t_1$, \dots,
  $t_{i-1}$, $t_{i+1}$, \dots, $t_n$ ही बंद पदे असतील, तर
\[\Atom{f}{t_1,\dots, t_{i-1}, t, t_{i+1}, \dots, t_n} \approx
\Atom{f}{t_1,\dots, t_{i-1}, t', t_{i+1}, \dots, t_n}.
\]
\item $t \approx t'$, $R$ हे विधेय, आणि $t_1$, \dots,
  $t_{i-1}$, $t_{i+1}$, \dots, $t_n$ ही बंद पदे असतील, तर
\begin{multline*}
\Atom{R}{t_1,\dots, t_{i-1}, t, t_{i+1}, \dots, t_n} \in \Gamma^* \text{ तेव्हा आणि केवळ तेव्हाच } \\
\Atom{R}{t_1,\dots, t_{i-1}, t', t_{i+1}, \dots, t_n} \in \Gamma^*.
\end{multline*}
\end{enumerate}
\end{prop}

\begin{proof}
$\Gamma^*$ सुसंगत आणि संपूर्ण असल्यामुळे $\eq[t][t'] \in
\Gamma^*$ तेव्हा आणि केवळ तेव्हाच, जेव्हा $\Gamma^* \Proves
\eq[t][t']$. म्हणून पुढील विधाने दाखवणे पुरेसे आहे:
\begin{enumerate}
\item प्रत्येक बंद पद~$t$ साठी $\Gamma^* \Proves \eq[t][t]$.
\item $\Gamma^* \Proves \eq[t][t']$ असेल, तर $\Gamma^* \Proves \eq[t'][t]$.
\item $\Gamma^* \Proves \eq[t][t']$ आणि $\Gamma^* \Proves
  \eq[t'][t'']$ असेल, तर $\Gamma^* \Proves \eq[t][t'']$.
\item $\Gamma^* \Proves \eq[t][t']$ असेल, तर
\[\Gamma^* \Proves
\eq[\Atom{f}{t_1,\dots,t_{i-1},t,t_{i+1},\dots,t_n}][\Atom{f}{t_1,\dots,t_{i-1},t',t_{i+1},\dots,t_n}]
\]
प्रत्येक $n$-स्थानी फलन~$f$ आणि बंद पदे $t_1$, \dots,
$t_{i-1}$, $t_{i+1}$, \dots,~$t_n$ यांच्यासाठी.
\item $\Gamma^* \Proves \eq[t][t']$ आणि
$\Gamma^* \Proves
\Atom{R}{t_1,\dots,t_{i-1},t,t_{i+1},\dots,t_n}$ असेल, तर
\[\Gamma^* \Proves
\Atom{R}{t_1,\dots,t_{i-1},t',t_{i+1},\dots,t_n}
\]
प्रत्येक $n$-स्थानी विधेय~$R$ आणि बंद पदे $t_1$, \dots,
$t_{i-1}$, $t_{i+1}$, \dots,~$t_n$ यांच्यासाठी.
\end{enumerate}

\end{proof}

\begin{prob}
\ref{fol:com:ide:prop:approx-equiv} ची सिद्धता पूर्ण करा.
\end{prob}

\begin{defn}
$\Gamma^*$ हा भाषा~$\Lang L$ मधील सुसंगत आणि संपूर्ण संच आहे,
$t$ हे बंद पद आहे, आणि~$\approx$ हा मागील व्याख्येतील संबंध आहे असे
समजा. मग:
\[\equivrep{t}{\approx} = \Setabs{t'}{t'\in \Trm[L], t \approx t'}
\]
आणि $\equivclass{\Trm[L]}{\approx} = \Setabs{\equivrep{t}{\approx}}{t \in \Trm[L]}$.
\end{defn}

\begin{defn}
\label{fol:com:ide:defn:term-model-factor}
$\Struct M = \Struct M(\Gamma^*)$ हे \ref{fol:com:mod:defn:termmodel} मधील
$\Gamma^*$ चे पद-प्रतिमान असू द्या. मग
$\Struct{\equivclass{M}{\approx}}$ ही पुढील संरचना आहे:
\begin{enumerate}
\item $\Domain{\equivclass{M}{\approx}} = \equivclass{\Trm[L]}{\approx}$.
\item $\Assign{c}{\equivclass{M}{\approx}} = \equivrep{c}{\approx}$
\item $\Assign{f}{\equivclass{M}{\approx}}(\equivrep{t_1}{\approx}, \dots,
  \equivrep{t_n}{\approx}) = \equivrep{\Atom{f}{t_1,\dots, t_n}}{\approx}$
\item $\tuple{\equivrep{t_1}{\approx}, \dots, \equivrep{t_n}{\approx}} \in
  \Assign{R}{\equivclass{M}{\approx}}$ तेव्हा आणि केवळ तेव्हाच, जेव्हा
  $\Sat{M}{\Atom{R}{t_1,\dots, t_n}}$; म्हणजे तेव्हा आणि केवळ तेव्हाच,
  जेव्हा $\Atom{R}{t_1,\dots, t_n} \in \Gamma^*$.
\end{enumerate}
\end{defn}

\begin{explain}
लक्षात घ्या की~$\equivclass{\Trm[L]}{\approx}$ च्या घटकांना
$\equivrep{t}{\approx}$ असे संबोधून, म्हणजेच
\emph{प्रतिनिधी}~$t \in \equivrep{t}{\approx}$ यांच्या साहाय्याने,
आपण त्यांच्यासाठी $\Assign{f}{\equivclass{M}{\approx}}$ आणि
$\Assign{R}{\equivclass{M}{\approx}}$ परिभाषित केली आहेत. या व्याख्या
प्रतिनिधींच्या निवडीवर अवलंबून नाहीत याची आपल्याला खात्री केली पाहिजे.
म्हणजे, तेच तुल्यतावर्ग निश्चित करणारी दुसरी~$t'$ निवडली
($\equivrep{t}{\approx} = \equivrep{t'}{\approx}$), तरी व्याख्यांतून तेच
फल मिळाले पाहिजे. उदाहरणार्थ, $R$ हे एकस्थानी विधेय असेल, तर
व्याख्येच्या शेवटच्या बाबीनुसार $\equivrep{t}{\approx} \in
\Assign{R}{\equivclass{M}{\approx}}$ तेव्हा आणि केवळ तेव्हाच, जेव्हा
$\Sat{M}{\Atom{R}{t}}$. आता $t \approx t'$ असलेल्या दुसऱ्या
पद~$t'$ साठी $\Sat/{M}{\Atom{R}{t'}}$ असेल, तर व्याख्येनुसार
$\equivrep{t'}{\approx} \notin \Assign{R}{\equivclass{M}{\approx}}$ असावे
लागेल. $t \approx t'$ असेल, तर $\equivrep{t}{\approx} =
\equivrep{t'}{\approx}$; पण $\equivrep{t}{\approx} \in
\Assign{R}{\equivclass{M}{\approx}}$ आणि $\equivrep{t}{\approx}
\notin \Assign{R}{\equivclass{M}{\approx}}$ ही दोन्ही विधाने लागू होऊ
शकत नाहीत. मात्र \ref{fol:com:ide:prop:approx-equiv} मुळे असे होऊ शकत नाही, याची
हमी मिळते.

\end{explain}

\begin{prop}
$\Struct{\equivclass{M}{\approx}}$ सुव्याख्यायित आहे; म्हणजे $t_1$,
  \dots, $t_n$, $t_1'$, \dots, $t_n'$ ही बंद पदे असतील आणि
  $t_i \approx t_i'$ असेल, तर
\begin{enumerate}
\item $\equivrep{\Atom{f}{t_1,\dots, t_n}}{\approx} =
    \equivrep{\Atom{f}{t_1',\dots, t_n'}}{\approx}$, म्हणजे
  \[  \Atom{f}{t_1,\dots, t_n} \approx \Atom{f}{t_1',\dots, t_n'}
  \]
  आणि
\item $\Sat{M}{\Atom{R}{t_1,\dots, t_n}}$ तेव्हा आणि केवळ तेव्हाच,
  जेव्हा $\Sat{M}{\Atom{R}{t_1',\dots, t_n'}}$; म्हणजे
  \[    \Atom{R}{t_1,\dots, t_n} \in \Gamma^* \text{ तेव्हा आणि केवळ तेव्हाच }
    \Atom{R}{t_1',\dots, t_n'} \in \Gamma^*.
  \]
\end{enumerate}
\end{prop}

\begin{proof}
\ref{fol:com:ide:prop:approx-equiv} वरून~$n$ वरील विगमनाने मिळते.
\end{proof}

पद-प्रतिमानाच्या बाबतीत होते त्याप्रमाणेच, सत्यता पूर्वप्रमेय सिद्ध
करण्यापूर्वी आपल्याला पुढील पूर्वप्रमेयाची गरज आहे.

\begin{lem}
\label{fol:com:ide:lem:val-in-termmodel-factored} $\Struct M = \Struct M
 (\Gamma^*)$ असू द्या; मग $\Value{t}{\equivclass{M}{\approx}} = \equivrep{t}
 {\approx}$.
\end{lem}
\begin{proof}
ही सिद्धता \ref{fol:com:mod:lem:val-in-termmodel} च्या सिद्धतेसारखीच आहे.
\end{proof}

\begin{prob}
\ref{fol:com:ide:lem:val-in-termmodel-factored} ची सिद्धता पूर्ण करा.
\end{prob}

\begin{lem}
\label{fol:com:ide:lem:truth}
सर्व वाक्ये~$\varphi$ यांच्यासाठी $\Sat{\equivclass{M}{\approx}}{\varphi}$ तेव्हा आणि
केवळ तेव्हाच, जेव्हा $\varphi \in \Gamma^*$.
\end{lem}

\begin{proof}
\ref{fol:com:mod:lem:truth} च्या सिद्धतेप्रमाणेच~$\varphi$ वरील विगमनाने सिद्धता
मिळते. $\varphi \ident \eq[t][t']$ असण्याच्या प्रसंगाकडेच फक्त अधिक लक्ष
देण्याची गरज आहे.
\begin{align*}
\Sat{\equivclass{M}{\approx}}{\eq[t][t']} & \text{ तेव्हा आणि केवळ तेव्हाच } \equivrep{t}{\approx} = \equivrep{t'}{\approx}
\text{ ($\Struct{\equivclass{M}{\approx}}$ च्या व्याख्येनुसार)}\\
& \text{ तेव्हा आणि केवळ तेव्हाच } t \approx t' \text{ ($\equivrep{t}{\approx}$ च्या व्याख्येनुसार)}\\
& \text{ तेव्हा आणि केवळ तेव्हाच } \eq[t][t'] \in \Gamma^* \text{ ($\approx$ च्या व्याख्येनुसार).}
\end{align*}
\end{proof}

\begin{digress}
लक्षात घ्या की $\Struct{M(\Gamma^*)}$ हे नेहमी गणनीय आणि अनंत
असले, तरी $\Struct{\equivclass{M}{\approx}}$ सांत असू शकते; कारण
$\equivrep{t}{\approx}$ असे फक्त सांत अनेक वर्ग असू शकतात. हे अपेक्षितच
आहे, कारण~$\Gamma$ मध्ये अशी वाक्ये असू शकतात की ती सत्य असणारी
कोणतीही संरचना सांत असणे आवश्यक असते. उदाहरणार्थ,
$\lforall[x][\lforall[y][x = y]]$ हे सुसंगत वाक्य आहे, पण ज्याच्या
प्रांत मध्ये नेमका एकच घटक आहे अशाच संरचनांमध्ये
त्याची पूर्ति होते.
\end{digress}












\section{संपूर्णता प्रमेय}\label{fol:com:cth:sec}

\begin{explain}
आपल्या निष्पत्ती एकत्र करू: आपल्याला संपूर्णता प्रमेय मिळते.
\end{explain}

\begin{thm}[संपूर्णता प्रमेय]
\label{fol:com:cth:thm:completeness}
$\Gamma$ हा वाक्यांचा संच असू द्या. $\Gamma$ सुसंगत असेल, तर
तो पूर्ततायोग्य आहे.
\end{thm}

\begin{proof}
$\Gamma$ सुसंगत आहे असे समजा.
  \ref{fol:com:hen:lem:henkin} नुसार असा संतृप्त सुसंगत संच
  $\Gamma' \supseteq \Gamma$ असतो. \ref{fol:com:lin:lem:lindenbaum} नुसार
असा $\Gamma^* \supseteq \Gamma'$ असतो की तो
सुसंगत आणि संपूर्ण असतो.

  $\Gamma' \subseteq \Gamma^*$ असल्यामुळे प्रत्येक
  सूत्र~$\varphi(x)$ साठी~$\Gamma^*$ मध्ये पुढील रूपाचे वाक्य
  असते:
  $\lexists[x][\varphi(x)] \lif \varphi(c)$.
  म्हणून~$\Gamma^*$ संतृप्त आहे. $\Gamma$ मध्ये~$\eq$ येत नसेल, तर
\ref{fol:com:mod:lem:truth} नुसार
$\Sat{M(\Gamma^*)}{\varphi}$ तेव्हा आणि केवळ तेव्हाच, जेव्हा $\varphi \in
\Gamma^*$. यावरून विशेषतः प्रत्येक $\varphi \in \Gamma$ साठी
$\Sat{M(\Gamma^*)}{\varphi}$; म्हणून
$\Gamma$ पूर्ततायोग्य आहे.
  $\Gamma$ मध्ये~$\eq$ येत असेल, तर \ref{fol:com:ide:lem:truth} नुसार
  सर्व वाक्ये~$\varphi$ यांच्यासाठी
  $\Sat{\equivclass{M}{\approx}}{\varphi}$ तेव्हा आणि केवळ तेव्हाच, जेव्हा
  $\varphi \in \Gamma^*$. विशेषतः, प्रत्येक $\varphi \in \Gamma$ साठी
  $\Sat{\equivclass{M}{\approx}}{\varphi}$; म्हणून $\Gamma$ पूर्ततायोग्य आहे.
\end{proof}

\begin{cor}[संपूर्णता प्रमेय, दुसरी आवृत्ती]
\label{fol:com:cth:cor:completeness}
प्रत्येक $\Gamma$ आणि प्रत्येक वाक्ये~$\varphi$ साठी: जर
$\Gamma \Entails \varphi$, तर $\Gamma \Proves \varphi$.
\end{cor}

\begin{proof}
लक्षात घ्या की \ref{fol:com:cth:cor:completeness} आणि
\ref{fol:com:cth:thm:completeness} यांमधील $\Gamma$ वर सर्ववाची संख्यकीकरण केलेले
आहे. आपला गोंधळ होऊ नये म्हणून \ref{fol:com:cth:thm:completeness} हे वेगळा चर
वापरून पुन्हा मांडू: वाक्यांचा कोणताही संच~$\Delta$ दिला असता,
जर $\Delta$ सुसंगत असेल, तर तो पूर्ततायोग्य आहे. प्रतिपरिवर्तनाने,
$\Delta$ पूर्ततायोग्य नसेल, तर $\Delta$ विसंगत आहे. अनुप्रमेय सिद्ध
करण्यासाठी याचा वापर करू.

$\Gamma \Entails \varphi$ असे समजा. मग \ref{fol:syn:sem:prop:entails-unsat}
नुसार $\Gamma \cup \{\lnot \varphi\}$ अपूर्ततायोग्य आहे. $\Gamma \cup
\{\lnot \varphi\}$ हाच आपला $\Delta$ घेतल्यास,
\ref{fol:com:cth:thm:completeness} च्या मागील आवृत्तीनुसार $\Gamma \cup
\{\lnot \varphi\}$ विसंगत आहे. पुढील संदर्भांनुसार,

  \ref{fol:axd:prv:prop:prov-incons}, \ref{fol:seq:prv:prop:prov-incons}, \ref{fol:ntd:prv:prop:prov-incons}, \ref{fol:tab:prv:prop:prov-incons},
$\Gamma \Proves \varphi$.
\end{proof}


\begin{prob}
\ref{fol:com:cth:cor:completeness} वापरून
\ref{fol:com:cth:thm:completeness} सिद्ध करा आणि अशा रीतीने
संपूर्णता प्रमेयाच्या दोन्ही मांडण्या समतुल्य आहेत हे दाखवा.
\end{prob}





\begin{prob}
एखादी निष्पत्ती प्रणाली संपूर्ण असण्यासाठी तिचे नियम इतके समर्थ
असले पाहिजेत की प्रत्येक अपूर्ततायोग्य संच विसंगत असल्याचे ते सिद्ध करू
शकतील. संपूर्णता सिद्ध करण्यासाठी निष्पत्तीचे कोणते नियम आवश्यक
होते? यांपैकी कोणताही नियम सिद्धतेत कोठेच वापरलेला नाही का? या प्रश्नांची
उत्तरे देण्यासाठी, \ref{fol:com:cth:thm:completeness} च्या सिद्धतेपर्यंत
नेणाऱ्या कोणत्या निष्पत्तींमध्ये निष्पत्तीचे कोणते नियम वापरले, हे
दाखवणारी यादी किंवा आकृती तयार करा. या सिद्धतांमधील नियमांच्या कोणत्याही
अध्याहृत वापरांची नोंद अवश्य करा.
\end{prob}















\section{संहतता प्रमेय}\label{fol:com:com:sec}

संपूर्णता प्रमेयाची एक महत्त्वाची निष्पत्ती म्हणजे संहतता प्रमेय.
वाक्यांच्या एखाद्या संचाचा प्रत्येक \emph{सांत} उपसंच पूर्ततायोग्य असेल,
तर संपूर्ण संच पूर्ततायोग्य असतो---तो संच स्वतः अनंत असला तरीही, असे
संहतता प्रमेय सांगते. हे मुळीच उघड नाही. प्रथमदर्शनी तरी अशी शक्यता
नाकारली जाईल असे काही दिसत नाही की वाक्यांचे अनंत संच व्याघाती
असतील, पण व्याघात, जणू काही, अनंत संख्येमुळेच उद्भवत असेल. अशी परिस्थिती
नाकारता येते, असे संहतता प्रमेय सांगते: ज्यांचा प्रत्येक सांत उपसंच
पूर्ततायोग्य आहे असे वाक्यांचे अपूर्ततायोग्य अनंत संच नसतात.
संपूर्णता प्रमेयाप्रमाणे याची तार्किक निष्पन्नतेशी संबंधित आवृत्तीही आहे:
वाक्यांच्या अनंत संचापासून काही निष्पन्न होत असेल, तर ते आधीच
एखाद्या सांत उपसंचापासून निष्पन्न होते.

\begin{defn}
  सूत्रांचा संच~$\Gamma$ \emph{सांततः पूर्ततायोग्य} असतो तेव्हा आणि
  केवळ तेव्हाच, जेव्हा प्रत्येक सांत $\Gamma_0 \subseteq \Gamma$
  पूर्ततायोग्य असतो.
\end{defn}

\begin{thm}[संहतता प्रमेय]
\label{fol:com:com:thm:compactness}
वाक्यांचा कोणताही संच~$\Gamma$ आणि कोणतेही वाक्य~$\varphi$ यांच्यासाठी पुढील
विधाने लागू होतात:
\begin{enumerate}
  \item $\Gamma \Entails \varphi$ तेव्हा आणि केवळ तेव्हाच, जेव्हा असा सांत
    $\Gamma_0 \subseteq \Gamma$ असतो की $\Gamma_0 \Entails \varphi$.
  \item $\Gamma$ पूर्ततायोग्य असतो तेव्हा आणि केवळ तेव्हाच, जेव्हा तो
    सांततः पूर्ततायोग्य असतो.
\end{enumerate}
\end{thm}

\begin{proof}
आपण (2) सिद्ध करू. $\Gamma$ पूर्ततायोग्य असेल, तर अशी
संरचना~$\Struct{M}$
असते की प्रत्येक $\varphi \in \Gamma$ साठी
$\Sat{M}{\varphi}$. अर्थात, ही
$\Struct{M}$ देखील~$\Gamma$ चा प्रत्येक सांत
उपसंच पूर्त करते; म्हणून~$\Gamma$ सांततः पूर्ततायोग्य आहे.

आता $\Gamma$ सांततः पूर्ततायोग्य आहे असे समजा. मग प्रत्येक सांत
उपसंच~$\Gamma_0 \subseteq \Gamma$ पूर्ततायोग्य आहे. निर्दोषतेनुसार
(
  \ref{fol:axd:sou:cor:consistency-soundness}, \ref{fol:seq:sou:cor:consistency-soundness}, \ref{fol:ntd:sou:cor:consistency-soundness}, \ref{fol:tab:sou:cor:consistency-soundness}),
प्रत्येक सांत उपसंच सुसंगत आहे. मग पुढील संदर्भांनुसार $\Gamma$ स्वतः
सुसंगत असला पाहिजे:

  \ref{fol:axd:ptn:prop:proves-compact}, \ref{fol:seq:ptn:prop:proves-compact}, \ref{fol:ntd:ptn:prop:proves-compact}, \ref{fol:tab:ptn:prop:proves-compact}.
संपूर्णतेनुसार (\ref{fol:com:cth:thm:completeness}), $\Gamma$ सुसंगत
असल्यामुळे तो पूर्ततायोग्य आहे.
\end{proof}


\begin{prob}
\ref{fol:com:com:thm:compactness} मधील (1) सिद्ध करा.
\end{prob}





\begin{ex}
उपपत्ती~$\Gamma$ च्या प्रत्येक प्रतिमान~$\Struct{M}$ मध्ये प्रत्येक
पद~$t$ हे अर्थातच~$\Domain{M}$ मधील घटक निर्देशित करते.
$\Domain{M}$ मधील प्रत्येक घटक देखील कोणत्यातरी पदाने निर्देशित
केला जातो, याची हमी देता येते का? दुसऱ्या शब्दांत, ज्यांची सर्व प्रतिमाने
बंद पदांनी आच्छादित आहेत अशा उपपत्ती~$\Gamma$ असतात का? $\Gamma$ ला
अनंत प्रतिमाने असतील, तर असे नसते, हे संहतता प्रमेय दाखवते. ते पुढीलप्रमाणे
दिसते: $\Struct{M}$ हे~$\Gamma$ चे अनंत प्रतिमान असू द्या आणि $c$ हे
$\Gamma$ च्या भाषेत नसलेले स्थिरांक असू द्या. $\Gamma$ ची
भाषा~$\Lang{L}$ मधील प्रत्येक पद~$t$ साठी असलेल्या सर्व
$\eq/[c][t]$ या वाक्यांचा~$\Delta$ हा संच असू द्या, म्हणजे
  \[  \Delta = \Setabs{\eq/[c][t]}{t \in \Trm[L]}.
  \]
$\Gamma \cup \Delta$ चा सांत उपसंच $\Gamma' \cup \Delta'$ असा लिहिता
येतो; येथे $\Gamma' \subseteq \Gamma$ आणि $\Delta' \subseteq \Delta$.
$\Delta'$ सांत असल्यामुळे त्यात फक्त सांत अनेक पदे असू शकतात.
त्यांपैकी कोणत्याही पदाने निर्देशित न केलेला $\Domain{M}$ मधील
घटक $a \in \Domain{M}$ असू द्या आणि $\Struct{M'}$ ही
$\Struct{M}$ सारखीच, पण $\Assign{c}{M'} = a$ देखील असलेली संरचना
असू द्या. $\Delta'$ मध्ये येणाऱ्या प्रत्येक~$t$ साठी
$a \neq \Value{t}{M}$ असल्यामुळे $\Sat{M'}{\Delta'}$. तसेच
$\Sat{M}{\Gamma}$, $\Gamma' \subseteq \Gamma$, आणि~$\Gamma$ मध्ये $c$
येत नसल्यामुळे $\Sat{M'}{\Gamma'}$ देखील. म्हणून $\Gamma \cup \Delta$
च्या प्रत्येक सांत उपसंच $\Gamma' \cup \Delta'$ साठी
$\Sat{M'}{\Gamma' \cup \Delta'}$. त्यामुळे $\Gamma \cup \Delta$ चा
प्रत्येक सांत उपसंच पूर्ततायोग्य आहे. संहततेनुसार $\Gamma \cup \Delta$
स्वतः पूर्ततायोग्य आहे. म्हणून $\Sat{M}{\Gamma \cup \Delta}$ अशी
प्रतिमाने असतात. असे प्रत्येक $\Struct{M}$ हे~$\Gamma$ चे प्रतिमान आहे,
पण ते बंद पदांनी आच्छादित नाही; कारण~$\Lang{L}$ मधील सर्व पदे~$t$
यांच्यासाठी $\Value{c}{M} \neq \Value{t}{M}$.
\end{ex}

\begin{ex}
विधेय~$<$, स्थिरांक~$\Obj{0}$, $\Obj{1}$, आणि
फलने~$+$, $\times$, व~$-$ असलेल्या भाषा~$\Lang{L}$ चा
विचार करा. प्रांत~$\Rat$ आणि उघड निर्वचनार्थ असलेल्या
संरचना~$\Struct{Q}$ मध्ये सत्य असलेल्या या भाषा मधील सर्व
वाक्यांचा संच~$\Gamma$ असू द्या. $\Gamma$ म्हणजे परिमेय
संख्यांविषयी सत्य असलेल्या~$\Lang{L}$ मधील सर्व वाक्यांचा संच.
अर्थात~$\Rat$ मध्ये (आणि~$\Real$ मध्येही) अशा संख्या~$r$ नसतात की त्या
$0$ पेक्षा मोठ्या, पण प्रत्येक $k \in \PosInt$ साठी $1/k$ पेक्षा लहान
असतात. अशी संख्या अस्तित्वात असती, तर ती \emph{अतिसूक्ष्म संख्या} असती:
शून्येतर, पण अनंतपटीने लहान. $\Gamma$ ची अतिसूक्ष्म संख्या असलेली प्रतिमाने
अस्तित्वात आहेत, हे दाखवण्यासाठी संहतता प्रमेय वापरता येते. आपल्या भाषेत
भागाकाराचे फलन नाही (शून्याने भागाकार अपरिभाषित असतो आणि
फलनांचे निर्वचन सर्वत्र परिभाषित फलनांनी केले पाहिजे). तरीही
$r < 1/k$ हे व्यक्त करता येते; कारण हे तेव्हा आणि केवळ तेव्हाच लागू होते,
जेव्हा $r \cdot k < 1$. आता $c$ हे नवीन स्थिरांक असू द्या आणि
$\Delta$ पुढीलप्रमाणे असू द्या:
\[\{\Obj{0} < c\}
\cup \Setabs{ c \times \num k < \Obj{1} }{k \in \PosInt}
\]
(येथे $\num{k} = (\Obj{1} + (\Obj{1} + \dots + (\Obj{1} +
\Obj{1})\dots))$ मध्ये $k$ इतके~$\Obj{1}$ आहेत). $\Delta$ च्या कोणत्याही
सांत उपसंच~$\Delta_0$ साठी असा~$K$ असतो की~$\Delta_0$ मधील सर्व
वाक्ये $c \times \num{k} < \Obj{1}$ यांमध्ये $k < K$ असते.
$\Assign{c}{Q'} = 1/K$ असे ठेवून $\Struct{Q}$ चा विस्तार
$\Struct{Q'}$ पर्यंत केला, तर कोणत्याही सांत $\Gamma_0 \subseteq \Gamma$
साठी $\Sat{Q'}{\Gamma_0 \cup \Delta_0}$, आणि म्हणून $\Gamma \cup \Delta$
सांततः पूर्ततायोग्य आहे (सराव: हे सविस्तर सिद्ध करा). संहततेनुसार
$\Gamma \cup \Delta$ पूर्ततायोग्य आहे. $\Gamma \cup \Delta$ च्या
कोणत्याही प्रतिमान~$\Struct{S}$ मध्ये एक अतिसूक्ष्म संख्या असते, ती म्हणजे
$\Assign{c}{S}$.
\end{ex}



\begin{prob}
अंकगणिताच्या प्रमाण प्रतिमान~$\Struct{N}$ मध्ये प्रत्येक सूत्र~$\num{n}
< x$ ची पूर्ति करणारा कोणताही घटक~$k \in \Domain{N}$ नाही
(येथे $\num{n}$ हे $n$ इतकी~$\prime$ चिन्हे असलेले
$\Obj{0}^{\prime\dots\prime}$ आहे). अंकगणिताच्या प्रमाण
प्रतिमान~$\Struct{N}$ मध्ये सत्य असलेली अंकगणिताच्या भाषेतील
वाक्ये ही प्रत्येक सूत्र~$\num{n} < x$ ची पूर्ति \emph{करणारा}
घटक असलेल्या संरचना~$\Struct{N'}$ मध्येही सत्य आहेत,
हे दाखवण्यासाठी संहतता प्रमेय वापरा.
\end{prob}



\begin{ex}
एकरूपता असलेले प्रथम-कोटी तर्कशास्त्र हे प्रांताच्या आकाराला
एखादी किमान मर्यादा असली पाहिजे, असे व्यक्त करू शकते हे आपल्याला माहीत
आहे: वाक्य~$\varphi_{\ge n}$ (जे ``किमान $n$ भिन्न वस्तू आहेत'' असे सांगते)
फक्त अशा रचनांमध्ये सत्य असते ज्यांत $\Domain{M}$ मध्ये किमान~$n$ वस्तू
असतात. म्हणून
  \[  \Delta = \Setabs{\varphi_{\ge n}}{n \ge 1}
  \]
असे घेतले, तर~$\Delta$ चे कोणतेही प्रतिमान अनंत असले पाहिजे. म्हणून
एखाद्या उपपत्तीला फक्त अनंत प्रतिमानेच आहेत याची हमी तिच्यात~$\Delta$
मिळवून देता येते: $\Gamma \cup \Delta$ ची प्रतिमाने म्हणजे~$\Gamma$ ची
सर्व आणि फक्त अनंत प्रतिमाने.

अशा रीतीने प्रथम-कोटी तर्कशास्त्र अनंतता व्यक्त करू शकते. पण ते सांतता
व्यक्त करू शकत नाही, हे संहतता प्रमेय दाखवते. कारण वाक्यांचा
एखादा संच~$\Lambda$ सर्व आणि फक्त सांत संरचनांमध्ये पूर्त होतो
असे समजा. मग $\Delta \cup \Lambda$ सांततः पूर्ततायोग्य आहे. का? सांत
$\Delta' \cup \Lambda' \subseteq \Delta \cup \Lambda$ असा संच असू
द्या; येथे $\Delta' \subseteq \Delta$ आणि $\Lambda' \subseteq \Lambda$.
$\varphi_{\ge n} \in \Delta'$ असेल अशा संख्यांपैकी सर्वांत मोठी संख्या~$n$
असू द्या. $\Lambda$ ची सर्व सांत रचनांमध्ये पूर्ति होत असल्यामुळे,
त्याला सांत अनेक, पण $\ge n$ इतके घटक असलेले
प्रतिमान~$\Struct{M}$ असते. पण मग $\Sat{M}{\Delta' \cup \Lambda'}$.
संहततेनुसार $\Delta \cup \Lambda$ ला अनंत प्रतिमान असते; पण हे
$\Lambda$ ची पूर्ति फक्त सांत संरचनांमध्ये होते या गृहीतकाच्या
विरुद्ध आहे.
\end{ex}













\section{संहतता प्रमेयाची प्रत्यक्ष सिद्धता}\label{fol:com:cpd:sec}

संपूर्णता प्रमेयाच्या सिद्धतेतील कल्पनाच वापरून, संपूर्णता प्रमेयाचा आधार
न घेता संहतता प्रमेय प्रत्यक्ष सिद्ध करता येते. संपूर्णता प्रमेयाच्या
सिद्धतेत आपण वाक्यांच्या सुसंगत संच~$\Gamma$ पासून सुरुवात केली,
त्याचा विस्तार करून वाक्यांचा सुसंगत, संतृप्त आणि
संपूर्ण संच~$\Gamma^*$ मिळवला; आणि मग~$\Gamma^*$ पासून रचलेल्या
पद-प्रतिमान~$\Struct{M(\Gamma^*)}$
मध्ये~$\Gamma$ मधील सर्व वाक्ये सत्य आहेत, म्हणजे~$\Gamma$
पूर्ततायोग्य आहे, हे दाखवले.

सांततः पूर्ततायोग्य वाक्यसंच पूर्ततायोग्य असतो हे दाखवण्यासाठी हीच पद्धत
वापरता येते. फक्त सत्यता पूर्वप्रमेयापर्यंत नेणाऱ्या निष्पत्तींच्या अशा
समांतर आवृत्त्या सिद्ध कराव्या लागतील की त्यांत ``सुसंगत'' ऐवजी ``सांततः
पूर्ततायोग्य'' वापरलेले असेल.

\begin{prop}
\label{fol:com:cpd:prop:fsat-ccs}
$\Gamma$ हा संपूर्ण आणि सांततः पूर्ततायोग्य आहे असे समजा. मग:
\begin{enumerate}
\item $(\varphi \land \psi) \in \Gamma$
  तेव्हा आणि केवळ तेव्हाच, जेव्हा $\varphi \in \Gamma$ आणि $\psi \in \Gamma$
  ही दोन्ही विधाने लागू होतात.

\item $(\varphi \lor \psi) \in \Gamma$ तेव्हा आणि केवळ तेव्हाच,
  जेव्हा $\varphi \in \Gamma$ किंवा $\psi \in \Gamma$ यांपैकी एक लागू होते.

\item $(\varphi \lif \psi) \in \Gamma$ तेव्हा आणि केवळ तेव्हाच,
  जेव्हा $\varphi \notin \Gamma$ किंवा $\psi \in \Gamma$ यांपैकी एक लागू होते.
\end{enumerate}
\end{prop}


\begin{prob}
\ref{fol:com:cpd:prop:fsat-ccs} सिद्ध करा. $\Proves$ चा वापर टाळा.
\end{prob}





\begin{lem}
\label{fol:com:cpd:lem:fsat-henkin} प्रत्येक सांततः पूर्ततायोग्य संच~$\Gamma$ चा
विस्तार करून संतृप्त आणि सांततः पूर्ततायोग्य संच~$\Gamma'$ मिळवता येतो.
\end{lem}



\begin{prob}
\ref{fol:com:cpd:lem:fsat-henkin} सिद्ध करा. (सूचना: महत्त्वाची
पायरी म्हणजे निष्पत्त्या किंवा सुसंगततेचा अजिबात आधार न घेता,
$\Gamma_n$ सांततः पूर्ततायोग्य असेल, तर $\Gamma_n \cup \{\theta_n\}$
देखील सांततः पूर्ततायोग्य आहे हे दाखवणे.)
\end{prob}



\begin{prop}\label{fol:com:cpd:prop:fsat-instances}
$\Gamma$ हा संपूर्ण, सांततः पूर्ततायोग्य आणि संतृप्त आहे असे समजा.
\begin{enumerate}
\item $\lexists[x][\varphi(x)] \in \Gamma$ तेव्हा आणि केवळ तेव्हाच,
  जेव्हा किमान एका बंद पद~$t$ साठी $\varphi(t) \in \Gamma$.
\item $\lforall[x][\varphi(x)] \in \Gamma$ तेव्हा आणि केवळ तेव्हाच,
  जेव्हा सर्व बंद पदे~$t$ यांच्यासाठी $\varphi(t) \in \Gamma$.
\end{enumerate}
\end{prop}



\begin{prob}
\ref{fol:com:cpd:prop:fsat-instances} सिद्ध करा.
\end{prob}


\begin{lem}
\label{fol:com:cpd:lem:fsat-lindenbaum} प्रत्येक सांततः पूर्ततायोग्य संच~$\Gamma$
चा विस्तार करून संपूर्ण आणि सांततः पूर्ततायोग्य संच~$\Gamma^*$
मिळवता येतो.
\end{lem}


\begin{prob}
\ref{fol:com:cpd:lem:fsat-lindenbaum} सिद्ध करा. (सूचना: महत्त्वाची
पायरी म्हणजे $\Gamma_n$ सांततः पूर्ततायोग्य असेल, तर $\Gamma_n \cup
\{\varphi_n\}$ किंवा $\Gamma_n \cup \{\lnot \varphi_n\}$ यांपैकी एक संच सांततः
पूर्ततायोग्य असतो हे दाखवणे.)
\end{prob}




\begin{thm}[संहतता]
\label{fol:com:cpd:thm:compactness-direct} $\Gamma$ पूर्ततायोग्य असतो तेव्हा आणि
केवळ तेव्हाच, जेव्हा तो सांततः पूर्ततायोग्य असतो.
\end{thm}

\begin{proof}
$\Gamma$ पूर्ततायोग्य असेल, तर अशी
संरचना~$\Struct{M}$
असते की प्रत्येक $\varphi \in \Gamma$ साठी
$\Sat{M}{\varphi}$. अर्थात, ही
$\Struct{M}$ देखील~$\Gamma$ चा प्रत्येक
सांत उपसंच पूर्त करते; म्हणून~$\Gamma$ सांततः पूर्ततायोग्य आहे.

आता $\Gamma$ सांततः पूर्ततायोग्य आहे असे समजा.
  \ref{fol:com:cpd:lem:fsat-henkin} नुसार असा सांततः पूर्ततायोग्य, संतृप्त संच
  $\Gamma' \supseteq \Gamma$ असतो. \ref{fol:com:cpd:lem:fsat-lindenbaum} नुसार
$\Gamma'$ चा विस्तार करून संपूर्ण आणि सांततः
पूर्ततायोग्य संच~$\Gamma^*$ मिळवता येतो, आणि $\Gamma^*$
  अजूनही संतृप्त असतो. \ref{fol:com:mod:defn:termmodel} मधीलप्रमाणे
पद-प्रतिमान~$\Struct{M(\Gamma^*)}$
रचा. लक्षात घ्या की \ref{fol:com:mod:prop:quant-termmodel} हे
$\Gamma^*$ सुसंगत आहे (किंवा त्या दृष्टीने संपूर्ण वा संतृप्त आहे)
या तथ्यावर अवलंबून नव्हते; ते फक्त $\Struct{M(\Gamma^*)}$ हे बंद पदांनी
आच्छादित आहे या तथ्यावर अवलंबून होते.
सत्यता पूर्वप्रमेयाची सिद्धता (\ref{fol:com:mod:lem:truth}) तशीच लागू होते,
जर \ref{fol:com:ccs:prop:ccs} च्या संदर्भाऐवजी
\ref{fol:com:cpd:prop:fsat-ccs} चा संदर्भ वापरला आणि
\ref{fol:com:hen:prop:saturated-instances} च्या संदर्भाऐवजी
\ref{fol:com:cpd:prop:fsat-instances} चा संदर्भ वापरला.

\end{proof}


\begin{prob}
\ref{fol:com:cpd:thm:compactness-direct} च्या सिद्धतेला आवश्यक असलेली
सत्यता पूर्वप्रमेयाची (\ref{fol:com:mod:lem:truth}) संपूर्ण सिद्धता
लिहा.
\end{prob}














\section{लोव्हेनहाइम--स्कोलेम प्रमेय}\label{fol:com:dls:sec}

एखाद्या उपपत्तीला अनंत प्रतिमान असेल, तर तिला फार तर गणनीय अनंत
असलेले प्रतिमानही असते, असे लोव्हेनहाइम--स्कोलेम प्रमेय सांगते. या
तथ्याची तात्काळ निष्पत्ती अशी की एखाद्या संरचनेचा आकार
अगणनीय आहे हे प्रथम-कोटी तर्कशास्त्र व्यक्त करू शकत नाही: सर्व
अगणनीय संरचनांमध्ये पूर्त होणारे कोणतेही वाक्य
किंवा वाक्यांचा कोणताही संच हा एखाद्या गणनीय रचनेतही
पूर्त होतो.

\begin{thm}
\label{fol:com:dls:thm:downward-ls} $\Gamma$ सुसंगत असेल, तर त्याला
गणनीय प्रतिमान असते; म्हणजे ज्याचा प्रांत सांत किंवा
गणनीय अनंत आहे अशा संरचनांमध्ये तो पूर्ततायोग्य असतो.
\end{thm}

\begin{proof}
$\Gamma$ सुसंगत असेल, तर संपूर्णता प्रमेयाच्या सिद्धतेतून मिळणाऱ्या
संरचना~$\Struct M$ चा प्रांत~$\Domain{M}$ हा भाषा~$\Lang L$ मधील
पदांच्या संचापेक्षा मोठा नसतो. म्हणून $\Struct M$ फार तर
गणनीय अनंत आहे.
\end{proof}

\begin{thm}
\label{fol:com:dls:noidentity-ls} $\Gamma$ हा एकरूपतेविना प्रथम-कोटी तर्कशास्त्राच्या
भाषेतील वाक्यांचा सुसंगत संच असेल, तर त्याला गणनीय अनंत
प्रतिमान असते; म्हणजे ज्याचा प्रांत अनंत आणि गणनीय आहे अशा
संरचनांमध्ये तो पूर्ततायोग्य असतो.
\end{thm}

\begin{proof}
$\Gamma$ सुसंगत असेल आणि त्यात एकरूपता येणारी वाक्ये नसतील, तर संपूर्णता
प्रमेयाच्या सिद्धतेतून मिळणाऱ्या संरचना~$\Struct M$ चा
प्रांत~$\Domain{M}$ हा भाषा~$\Lang L'$ मधील पदांच्या संचाशी एकरूप असतो.
म्हणून $\Trm[L']$ हे गणनीय अनंत असल्यामुळे $\Struct{M}$ देखील
गणनीय अनंत आहे.
\end{proof}

\begin{ex}[स्कोलेमची विरोधापत्ती]
त्सेर्मेलो--फ्रेंकेल संचसिद्धान्त~$\Log{ZFC}$ ही एक अतिशय समर्थ चौकट
आहे; संचांच्या आकारांविषयीच्या तथ्यांसह जवळजवळ सर्व गणिती विधाने तिच्यात
व्यक्त करता येतात. म्हणून, उदाहरणार्थ, वास्तव संख्यांचा संच~$\Real$ हा
अगणनीय आहे हे~$\Log{ZFC}$ सिद्ध करू शकते; कोणत्याही संचाचा
घातसंच त्या संचापेक्षा मोठा असतो हे कँटर यांचे प्रमेयही ती सिद्ध करू शकते;
आणि असेच पुढे. $\Log{ZFC}$ सुसंगत असेल, तर तिची सर्व प्रतिमाने अनंत
असतात. शिवाय, त्या सर्व प्रतिमानांत असे घटक असतात की उपपत्तीच्या
मते ते अगणनीय असतात; उदाहरणार्थ, नैसर्गिक संख्यांचा घातसंच
अस्तित्वात आहे हे~$\Log{ZFC}$ चे प्रमेय सत्य करणारा घटक. लोव्हेनहाइम--स्कोलेम
प्रमेयानुसार $\Log{ZFC}$ ला गणनीय प्रतिमानेही असतात---ज्यांत
``अगणनीय'' संच असतात, पण जी स्वतः गणनीय असतात.
\end{ex}



